% concatenated from main.tex
\documentclass[reqno]{amsart}
\usepackage{graphicx}
\usepackage{epsfig}
\usepackage{amssymb}
\usepackage{xcolor}
%\usepackage[margin=1in]{geometry}
\setlength{\textwidth}{4.98in}
\setlength{\textheight}{7.61in}
%\usepackage{fullpage}
\usepackage[
backend=biber,
style=alphabetic,
sorting=nyt
]{biblatex}
\addbibresource{References.bib}
\usepackage{amsfonts}
\usepackage{dsfont}
\usepackage{amsmath}
\usepackage{stmaryrd}
\usepackage{hyperref}
\usepackage{subfigure, verbatim}
\pagestyle{myheadings}
\setlength{\textwidth} {13cm}
\newtheorem{conjecture}{Conjecture}
\newtheorem{theorem}{Theorem}
\newtheorem*{theorem*}{Theorem}
\newtheorem{lemma}{Lemma}
\newtheorem{proposition}{Proposition}[section]
\theoremstyle{definition}
\newtheorem{definition}{Key Proposition}
\newtheorem{example}{Example}
\newtheorem{corollary}{Corollary}
%\newtheorem*{lemma2}{Lemma 2}
\newtheorem{xca}[theorem]{Exercise}
\theoremstyle{remark}
  \newtheorem{rem}{Remark}%%[section]
  \newtheorem*{nota}{Notation}
\numberwithin{equation}{section}
\newcommand\numberthis{\addtocounter{equation}{1}\tag{\theequation}}
\theoremstyle{remark}
\newtheorem{remark}[theorem]{Remark}
\newcommand{\abs}[1]{\lvert#1\rvert}
\newcommand{\e}{\'{e}}
\newcommand{\Prob}{\mathbb{P}}
\newcommand{\E}{\mathbb{E}}
\newcommand{\half}{\frac{1}{2}}
\newcommand{\g}{\mathcal{G}}
\newcommand{\f}{\mathcal{F}}
\newcommand{\me}{\mathcal{E}}
\newcommand{\ome}{\overline{\mathcal{E}}}
\newcommand{\bm}{Brownian motion }
\newcommand{\bmn}{Brownian motion}
\newcommand{\beq}{\begin{equation}}
\newcommand{\eeq}{\end{equation}}
\newcommand{\beqn}{\begin{equation*}}
\newcommand{\eeqn}{\end{equation*}}
\newcommand{\bea}{\begin{eqnarray}}
\newcommand{\eea}{\end{eqnarray}}
\newcommand{\bean}{\begin{eqnarray*}}
\newcommand{\eean}{\end{eqnarray*}}
\newcommand{\blankbox}[2]{  \parbox{\columnwidth}{\centering
    \setlength{\fboxsep}{0pt}    \fbox{\raisebox{0pt}[#2]{\hspace{#1}}}  }}
%\input{tcilatex}
\renewcommand{\labelenumi}{(\alph{enumi})}
\renewcommand{\labelenumii}{(\roman{enumii})}

\newcommand{\be}{\begin{enumerate}}
\newcommand{\ee}{\end{enumerate}}
\newcommand{\bi}{\begin{itemize}}
\newcommand{\ei}{\end{itemize}}
\newcommand{\bd}{\begin{description}}
\newcommand{\ed}{\end{description}}
\newcommand{\ds}{\displaystyle}
\newcommand{\vol}{\operatorname{vol}}
\newcommand{\bR}{\mathbf{R}}
\newcommand{\bN}{\mathbf{N}}
\newcommand{\bC}{\mathbf{C}}
\newcommand{\bZ}{\mathbf{Z}}
\newcommand{\bH}{\mathbf{H}}
\newcommand{\bP}{\mathbf{P}}
\newcommand{\cB}{\mathcal{B}}
\newcommand{\cS}{\mathcal{S}}
\newcommand{\cF}{\mathcal{F}}
\newcommand{\cH}{\mathcal{H}}
\newcommand{\cD}{\mathcal{D}}
%\newcommand{\e}{\varepsilon_\theta}
\newcommand{\I}{\operatorname{Id}}
\newcommand{\Hess}{\operatorname{Hess}}
\newcommand{\dd}{\mathrm{d}}
\newcommand{\pa}{\partial}
\newcommand{\wt}{\widetilde}
\newcommand{\la}{\langle}
\newcommand{\ra}{\rangle}

\begin{document}
\title[Critical partition function in the $q$-aspect]{On the $\beta=2$ Partition function for Dirichlet $L$-functions in the $q$-aspect}
\author{Christopher Atherfold}
\address{School of Mathematics, University of Bristol, Bristol, BS8 1UG, United Kingdom}
\email{christopher.atherfold@bristol.ac.uk}
\begin{abstract}
    We study the $\beta=2$ partition function $\int_{|h| \leq \log^{\theta}(q)/2}|L(1/2+ih,\chi)|^2dh$ for typical Dirichlet characters modulo a large prime $q$ and $\theta \in (-1/2,0]$ motivated by a $q$-analogue of the Saksman--Webb conjectures. When $\theta <0$, we use Harper's randomisation argument to introduce explicit conditioning to recover moment upper bounds consistent with critical normalisation predicted there. As an application, we prove that for $q(1-o(1))$ Dirichlet characters modulo $q$, $\max_{|h| \leq 1/2}|L(1/2+ih,\chi)|\ll \frac{\log(q)}{(\log\log(q))^{3/4+o(1)}}$, establishing an upper bound matching the predictions of the $q$-analogue of the Fyodorov--Hiary--Keating conjectures up to second order.
\end{abstract}
\maketitle
\section{Introduction} \subsection{Motivation and results}
We study how large $L$-functions can be for typical Dirichlet characters modulo a large prime $q$ near the central point through the lens of the $\beta=2$ partition function. This is given by
\begin{equation}
    \mathcal{Z}_2(q,\chi,\theta)=\int_{|h| \leq \log^{\theta}(q)/2}|L(1/2+ih,\chi)|^2dh
\end{equation}
where $\theta \in (-1/2,0]$. \par
Two conjectures which have shaped recent study of the Riemann zeta function are the Fyodorov--Hiary--Keating and Saksman--Webb conjectures, both of which we will discuss in this paper. \par
We begin with the Saksman--Webb conjectures. In \cite{sw20}, Saksman and Webb investigate many distributional questions concerning random Euler products which they conjecture should also hold for the Riemann zeta function on typical intervals. Among their conjectures is 
\begin{equation} \label{sw_zeta}
    \frac{\sqrt{\log\log(T)}}{\log(T)}\int_0^1|\zeta(1/2+i(\tau+h)|^2dh \underset{T \rightarrow \infty}{\longrightarrow} \int_0^1g(h)\lambda_2(dh)
\end{equation}
in distribution, where $\tau$ is a uniformly distributed random variable on the interval $[T,2T]$ and $\lambda_2(dh)$ is the critical Gaussian multiplicative chaos measure. The function $g$ is a generalised random function which is well behaved in various senses, and for more details one can see the theorem statements in \cite{sw20}. This conjecture will be settled in forthcoming work of Harper, Saksman and Webb. It is natural to ask if the analogue of such a conjecture holds for unitary averages over Dirichlet characters. We present the $q$-analogue of Conjecture $1.11$ in \cite{sw20}.
\begin{conjecture} \label{q-aspect_partition_conjecture}
Let $q$ be a large prime and $\chi_q$ to be a non-trivial Dirichlet character modulo $q$ chosen uniformly at random. Then
    $$\frac{\sqrt{\log\log(q)}}{\log(q)}\int_0^1|L(1/2+ih,\chi_q)|^2dh \underset{q \rightarrow \infty}{\longrightarrow} \int_0^1g(h)\lambda_2(dh),$$
    where $g(h)$ and $\lambda_2(dh)$ are the same as identified in equation (\ref{sw_zeta}).
\end{conjecture}
As mentioned in \cite{harper2019partition}, this can be described as the $\beta=2$ partition function of $\log(|L(1/2+ih,\chi)|)$ for fixed $\chi$ and varying $h$. Two striking features of the conjecture are the normalisation factor and the non-Gaussian limit. Before determining the limiting distribution, we need to have the correct normalisation. Our first result provides evidence towards this being the correct normalisation factor, together with the work of Harper in \cite{harper2023typicalsizecharacterzeta} and the recent work of Vihko \cite{vihko2025dirichlet} where Vihko proves that a uniformly random Dirichlet $L$-function on the critical line converges as a random Schwartz distribution to the Steinhaus random Euler product as $q \to \infty$ is a positive integer. Our proofs use the randomisation argument developed by Harper in \cite{harper2019partition} and \cite{harper2023typicalsizecharacterzeta} to exploit the connection between the deterministic quantities and statements involving Steinhaus random multiplicative functions. The results in \cite{harper2023typicalsizecharacterzeta} are particularly revealing, since Harper proves that for $q$ a large prime and $x=o(q)$ then
$$\frac{1}{q-1}\sum_{\chi \bmod{q}}\big|\sum_{n \leq x}\chi(n)\big| = o(\sqrt{x}).$$
In the setting of Conjecture \ref{q-aspect_partition_conjecture}, we obtain the following theorem.
\begin{theorem}
Suppose $q$ is a large prime. Then uniformly for $r \in [0,1]$,
    $$\frac{1}{q-1}\sum_{\chi \bmod{q}}(\mathcal{Z}_2(q,\chi,0))^r \ll \big(\frac{\log(q)}{1+(1-r)\sqrt{\log\log(q)}}\big)^r.$$
    \label{pf_moments_constant}
\end{theorem}
We also study the $\beta=2$ partition function on mesoscopic (shrinking) intervals of size $\log^{\theta}(q)$ for $\theta \in (-1/2,0]$. This introduces a new difficulty, where we need to introduce explicit conditioning on the contribution of $P_\theta = \exp(\log^{|\theta|}(q))$ smooth integers. We have for $P_\theta$ smooth integers, $n^{ih}$ is essentially frozen up to $O(1)$ fluctuations. Consequently, we need to condition on the value of the Dirichlet polynomial over the $P_\theta$-smooth numbers to obtain the correct order of magnitude. This has previously been studied in Arguin--Hamdan \cite{ah24} and before that in \cite{aor21}. Our approach to this is rather different. Previously, the authors used level sets of $\log(|\zeta|)$ to derive their results, whereas we are integrating $L$ directly, which necessitates the explicit conditioning. \par
In \cite{harper2023typicalsizecharacterzeta}, Harper can implement the \textit{coarser conditioning} implicitly, whereas we must do this explicitly. We achieve this by reformulating the multipoint implicit conditioning there into an (explicit) one point conditioning statement using a chaining argument. Suppose $f(n)$ be a Steinhaus random multiplicative function, so define the sequence $(f(p))$ to be independent identically distributed random variables on the complex unit circle and extended completely multiplicatively to the natural numbers. Take
\begin{equation} \label{short_lognormal}
        Y_x(h) := \Re \big(\sum_{p \leq x}\big(\frac{f(p)}{p^{1/2+ih}}+\frac{(f(p))^2}{p^{1+2ih}}\big)\big),
\end{equation}
and throughout the paper, we assume $W$ is some real number.
\begin{theorem}\label{pf_moments_meso}
    Fix $\theta \in (-1/2,0]$ and let $P_\theta = \exp(\log^{|\theta|}(q))$. Then, uniformly for $q$ a large prime and $r \in [0,1]$, if $\theta <0$
    \begin{equation*}
        \frac{1}{q-1}\sum_{\chi \bmod{q}}\big[(\mathcal{Z}_2(q,\chi,\theta))^r\big|Y_{P_\theta}(0) \in [W,W+1]\big] \ll \big(\frac{e^{2W}\log^{1+2\theta}(q)}{1+(1-r)\sqrt{\log\log(q)}}\big)^r.
    \end{equation*}
    We also have for $W \in \mathbb{R}$
    \begin{equation*}
        \frac{1}{q-1}\sum_{\chi \bmod{q}}\big[(\mathcal{Z}_2(q,\chi,\theta)^r\big|Y_{P_\theta}(0) \leq W\big] \ll \big(\frac{e^{2W}\log^{1+2\theta}(q)}{1+(1-r)\sqrt{\log\log(q)}}\big)^r.
    \end{equation*}
\end{theorem}
The second estimate immediately follows from the first splitting the interval $(-\infty,W]$ into subintervals of unit length and applying the first statement in Theorem \ref{pf_moments_meso} to each subinterval. The conditioning in the second statement is much closer to a barrier event on the frozen field, which are common in the context of log-correlated fields and branching random walks. As a corollary of Theorems \ref{pf_moments_constant} and \ref{pf_moments_meso}, we obtain the following tail bounds.
\begin{corollary} \label{typical_partition_size}
    \textit{Fix $\theta \in (-1/2,0)$. Then, uniformly for a large prime $q$ and $\lambda \geq 2$ if $\theta <0$}
     \begin{align*}
        \frac{1}{q-1}\#\big\{\chi \bmod{q}:\mathcal{Z}_2(q,\chi,\theta) > \frac{\lambda e^{2W} \log^{1+2\theta}(q)}{\sqrt{\log\log(q)}} \big| Y_{P_\theta}(0) \leq W\big\} 
        \ll \frac{\min\{\log(\lambda),\sqrt{\log\log(q)}\}}{\lambda}
    \end{align*}
    \textit{For $\theta=0$, we have}
    \begin{align*}
        \frac{1}{q-1}\#\big\{\chi \bmod{q}:\mathcal{Z}_2(q,\chi,\theta) > \frac{\lambda \log(q)}{\sqrt{\log\log(q)}}\big\} 
        \ll \frac{\min\{\log(\lambda),\sqrt{\log\log(q)}\}}{\lambda}.
    \end{align*}
\end{corollary}
\begin{proof}
    When $\lambda \geq \exp(\sqrt{\log\log(q)})$, then this is an application of Markov's inequality when $r=1$ with the moment bound obtained in Theorems \ref{pf_moments_constant} and \ref{pf_moments_meso}. Otherwise, one applies Markov's inequality to the power $r=1-\frac{2}{\log(\lambda)}$.
\end{proof}
The corollary indicates that the normalisation in Conjecture \ref{q-aspect_partition_conjecture} is the correct order of magnitude. The mesoscopic theorem can also be generalised to the $t$-aspect for the Riemann zeta function as well, which improves upon bounds obtained in \cite{aor21} and \cite{ah24} in our range of $\theta$. \par
The appearance of the $\sqrt{\log\log(q)}$ in the theorems is connected to the exponent $\beta=2$, the critical exponent in the language of multiplicative chaos. At the critical exponent, there are no longer enough sufficiently high points to support a non-degenerate measure, so we require the Seneta--Heyde normalisation to obtain something meaningful, which corresponds to the $\sqrt{\log\log(q)}$ correction. From a number theoretic perspective, we expect the characters and points which contribute to the size of $\sum_{\chi \bmod{q}}\int_{|h| \leq 1/2}|L(1/2,\chi)|^2dh$ to satisfy $|L(1/2+ih,\chi)| \approx \log(q)$. Such pairs $(\chi,h)$ are rare when one considers $h$ over a constant length or shrinking interval, which necessitates the renormalisation. This is described in greater detail in \cite{sound09}. In very broad terms, this is why the $\beta=2$ partition function is so important for understanding the typical maxima of these $L$-functions. This type of correction is also seen when the intervals considered are growing very slowly. This is investigated for a random model of zeta in the work of Chang \cite{chang2024}, and also appears when studying sums of random multiplicative functions over short intervals. For more details on the latter, see Caich \cite{c24} and Harper--Soundararajan--Xu \cite{hsx26}. \par
The Fyodorov--Hiary--Keating conjectures were first proposed in \cite{fhk12} and explored further in \cite{fk14}. These have been studied rigourously in many works since these papers, with state of the art due to Arguin--Bourgade--Radziwi\l{}\l{} in \cite{abr20} and \cite{abr23}. Among other results, they prove that for $T$ large and $\tau$ being a uniform random variable taking values in $[T,2T]$, then
$$\mathbb{P}_\tau\big(\max_{|h| \leq 1}|\zeta(1/2+i(\tau+h))| > \frac{e^U\log(T)}{(\log\log(T))^{3/4}}\big) \asymp Ue^{-2U}$$
for $U \in [10,\sqrt{\log\log(T)}]$. This confirms order of magnitude for the size of the right tail decay. As for the latest developments in the random matrix setting, see the work of Paquette and Zeitouni \cite{pz22}. We propose an equivalent statement for the maxima on the unit interval centred at $s=1/2$ for typical Dirichlet characters modulo a large prime $q$.
\begin{conjecture} \label{constant_qFHK}
    Suppose $q$ is a sufficiently large prime. Then one has that for $0 \leq U \leq \log\log(q)$,
    \begin{equation} \label{q_count}
        \frac{1}{q-1}\#\big\{\chi \bmod{q}: \max_{|h| \leq 1/2}|L(1/2+ih,\chi)| \geq \frac{e^U\log(q)}{(\log\log(q))^{3/4}} \big\} \asymp Ue^{-2U}.
    \end{equation}
\end{conjecture}
While the work of Harper in \cite{harper2023typicalsizecharacterzeta} predicts the normalisation in Conjecture \ref{q-aspect_partition_conjecture}, it is unclear from previous literature that such a conjecture should hold as the typical maxima are significantly more delicate. In our final theorem, we prove that this conjecture is sharp in the upper bound up to the second order correction term.
\begin{theorem}
    Suppose $q$ is a large prime, $0 \leq U \leq \log\log(q)$. Then 
    \begin{align*}
        &\frac{1}{q-1}\#\big\{\chi \bmod{q}: \max_{|h| \leq 1/2}|L(1/2+ih,\chi)| \geq \frac{e^U\log(q)}{(\log\log(q))^{3/4}} \big\} \\ \ll& e^{-2U}(\log\log\log(q))^2(\log\log\log(q) + U).
    \end{align*}
    \label{short_range_max}
\end{theorem}
In particular, for any $U(q) \to \infty$, we have
$$\#\big\{\chi \bmod{q}: \max_{|h| \leq 1/2}|L(1/2+ih,\chi)| \geq \frac{e^U\log(q)(\log\log\log(q))^{3/2}}{(\log\log(q))^{3/4}}\big\} =o(q).$$
\subsection{Proof ideas}
There are two main components to each proof: a randomisation argument and a probabilistic analysis of the resulting random Euler products. In all our theorems, we adapt Harper's approach in \cite{harper2019partition} to our setting. \par
For all three of our theorems, once we transfer to sums of random multiplicative functions, the probabilistic analysis is rather similar as we can transition to studying expressions of the form
\begin{equation}\label{heuristic_quantity}
    \sum_{\chi \bmod{q}}\mathbf{1}_{\mathcal{G}_\chi}\int_{|h| \leq 1/2}|\sum_{\substack{m \leq q^{\varepsilon} \\ P^+(m) \leq P}}\frac{\chi(m)}{m^{1/2+ih}}\sum_{\substack{n \leq q^{1/2-2\varepsilon} \\ P^-(n) > P}}\frac{\chi(n)}{n^{\sigma+ih}}|^2dh
\end{equation}
where $P$ is a suitable parameter ($P$ is chosen such that $\log\log(P) \asymp \log\log(q)$) and $\mathcal{G}_\chi$ is an event which occurs with high probability which controls the size of the Dirichlet polynomial at many scales. Indeed, the key observation is after passing to the random side, the size of the Dirichlet polynomial over the $P$ smooth integers conditioned on the appropriate randomised event $\mathbf{1}_{\mathcal{G}_\text{rand}}$ is $\ll \log(P)/\sqrt{\log\log(P)}$. This is a consequence of the ballot theorem which is sufficient for the bounds in Theorems \ref{pf_moments_constant} and \ref{pf_moments_meso}.\par
The further saving of $\log\log(P)$ seen in Theorem \ref{short_range_max} comes from strengthening the event $\mathcal{G}_\chi$ to also include a condition that ensures that $$|\sum_{\substack{m \leq q^{\varepsilon} \\ P^+(m) \leq P}}\frac{\chi(m)}{m^{1/2+ih}}| > \frac{\log(P)}{V}$$
for some parameter $V$. Since we are looking at very large values, one hopes to save even more from the ballot theorem. This is the case, and we conclude from there. This demonstrates the importance of the $\beta=2$ partition function since it allows for the full range of ballot theorem savings to be accessed, which are simply not available at smaller $\beta$. \par
When setting up the transfer to random multiplicative functions, we use both the approximate functional equation for the $L$-function and then approximate this using a product of two Dirichlet polynomials (three in the case of Theorem \ref{pf_moments_meso}) split according to their prime factor compositions. At this stage, one is ready to apply Harper's randomisation estimates in \cite{harper2023typicalsizecharacterzeta} in the proofs of Theorem \ref{pf_moments_constant} and \ref{pf_moments_meso}. The orthogonality of Dirichlet characters disappears when we attempt to condition upon them directly, so we have to move to the random side so we can condition without affecting orthogonality. \par
The proof of Theorem \ref{short_range_max} requires a different initial approach. A similar idea applies, but we use the Cauchy integral formula around a typical maximum. This is done by approximating its value by an integral over a rectangle containing the maximum using the approximate functional equation
$$\sum_{n \leq \sqrt{q/2\pi}}\frac{\chi(n)}{n^{1/2+ih(*)}} = \frac{1}{2\pi i}\int_{\mathcal{R}(*)}\frac{\sum_{n \leq \sqrt{q/2\pi}}\frac{\chi(n)}{n^s}}{(s-1/2-ih(*))}ds$$
where $\mathcal{R}(*)$ is a small rectangle in the complex plane containing the point $h(*)$ (with sides of length $\asymp 1/\log(q)$. Eventually, it is sufficient to use the equivalent of a union bound over every point in the constant length interval as these maxima are very rare. Various applications of Cauchy-Schwarz and Holder's inequality allow us to reduce this to studying either (simple) fourth moment estimates or the $\beta=2$ partition function. \par
The transfer mechanism for the proof of Theorem \ref{pf_moments_meso} differs in one key aspect. The style of mollifier we use here is derived from mean value estimates of Dirichlet polynomials as we are approximating $L$ directly (as opposed to $\log(|L|)$ as seen in \cite{abr20} for example). Additive mollifiers may be conceptually simpler to work with, but they are not capable of providing a good enough approximation if $\theta$ is too small. We see this when $\theta$ approaches $-1/2$ suitably quickly, then the mean value estimates are not sufficient for obtaining the correct order of magnitude. \par 
Finally, we mention the chaining argument we use to extract the mesoscopic conditioning factor. This follows similarly to Arguin--Hamdan in \cite{ah24}, and we get to that point once we successfully transfer from character sums over smooth numbers to sums of random multiplicative functions over smooth numbers. We do this to avoid periodicity issues with Dirichlet characters which the Steinhaus random multiplicative functions do not have. When on the random side, we can approximate this Dirichlet polynomial using (something very close to) a truncated Euler product. A similar strategy is seen in the work of Hardy and Xu in \cite{hx26}.
\subsection{Further discussion of results} \label{further_disc}
Given our discussion of Theorems \ref{pf_moments_constant} and \ref{pf_moments_meso}, we remark on Conjecture \ref{constant_qFHK} and Theorem \ref{short_range_max}. Following the work of Arguin--Bourgade--Radziwi\l{}\l{}, it should be possible to recover Conjecture \ref{constant_qFHK} in its entirety if one adapts the work in \cite{abr20} and \cite{abr23} to the $q$-aspect. Since we focus on an upper bound in this paper, the arithmetic input in \cite{abr20} could be mimicked by computing a similar twisted fourth moment in the $q$-aspect (such an argument would be similar to that seen in \cite{ac25} and \cite{creighton25} with the fourth moment instead of the second moment they study there). However, this is typically a significant algebraic challenge (see the work of \cite{gz25twisted} for example) to obtain a tractable expression, let alone obtaining bounds in $q$. We still achieve the second order correction term in Theorem \ref{short_range_max}, and could be further improved by a more delicate handling of the Dirichlet polynomial over the rough numbers. \par
As for improvements to the results, the most significant one is to Theorem \ref{short_range_max} since we resort to approximating a Dirichlet polynomial over rough numbers by Cauchy--Schwarz, which results in forcing one of the parameters $V$ to be taken as too large later on in the proof. The other issue is a slight loss on the sharpness of the right tail upper bound due to Markov's inequality in Corollary \ref{typical_partition_size}. We conjecture that for $\theta \in (-1/2,0]$ and fixed $\lambda \geq 2$, then one has
$$\frac{1}{q-1}\# \big\{\chi \bmod{q}: \mathcal{Z}_2(q,\chi,\theta) > \frac{\lambda \log(q)}{\sqrt{(1+\theta)\log\log(q)}}\big\} \asymp \frac{1}{\lambda}$$
in line with similar types of results obtained in \cite{barral15}, which would likely require Conjecture \ref{q-aspect_partition_conjecture} to be proved fully. \par
Finally, we discuss matching lower bounds for our theorems. For Theorem \ref{short_range_max}, following the work of \cite{abr23} could lead to a lower bound at the precision predicted by Conjecture \ref{constant_qFHK} as their approach is very probabilistic in nature. However, proving an order of magnitude result comparable to Theorem \ref{pf_moments_constant} is a different challenge. This would either require a proof of Conjecture \ref{q-aspect_partition_conjecture}, or deploying a strategy similar to that in \cite{harper26}. These need rather different techniques to the ones developed here, so we choose not to pursue them.
\subsection{Notation and organisation of the paper}
Due to the different notions of averaging throughout the paper, we summarise them here. We will use $\mathbb{E}^{\text{char}}$ to denote averaging over all Dirichlet characters modulo $q$, and simply $\mathbb{E}$ to refer to averages involving random multiplicative functions, and when we want to condition on the value of short Dirichlet polynomial $Y_{P_\theta}(0)$ to be in the interval $[W,W+1]$ we denote this by $\mathbb{E}_{W}^{\text{char}}$ and $\mathbb{E}_W$ respectively. There are also times when we want to understand the size of a set where an event occurs for characters, so we let $\mathbb{P}_q$ denote the uniform probability measure over the set of Dirichlet characters mod $q$. Since we are also considering the mesoscopic regime, we fix some important quantities. For a given value of $\theta$, we let $P_\theta = \exp(\log^{|\theta|}(q))$ and $\varepsilon_\theta = \frac{\log^\theta(q)}{(\log\log(q))^2}$ (which we observe is $\gg (\log(q))^{-1/2}$ for the range of $\theta$ we are investigating). Finally we discuss Dirichlet polynomials where the contributions come from integers with certain factorisation properties throughout the paper. We use $P^+(n)$ and $P^{-}(n)$ to denote the largest and smallest prime factor of $n$ respectively.\par
We collect most of our results we need to prove our theorems in Section \ref{sec2}. We then prove the chaining argument we use throughout in Section \ref{chain_arg}. Theorems \ref{pf_moments_constant} and \ref{pf_moments_meso} will be proven in Sections \ref{pfm_proof} and both theorems will be proved together. We defer the proof of Theorem \ref{short_range_max} to Section \ref{typical_max_section}. Where appropriate, we will also comment on the differences in approaches seen in \cite{harper2019partition} and \cite{abr20}. For a more general review of results relating to the progress on the FHK conjectures and critical multiplicative chaos, see the papers of Bailey--Keating \cite{bk22} and Powell \cite{powell22} respectively.
\subsection*{Acknowledgements}
The author would like to thank his supervisors Oleksiy Klurman and Joseph Najnudel for carefully reading an earlier draft of this paper, and for various discussions relating to this work. He also thanks Jad Hamdan in particular for sharing a version of \cite{ah24} early and for various related discussions. The author also thanks Emma Bailey and Andrew Pearce-Crump for their comments and encouragement on this problem. Finally, he thanks Adam Harper, Seth Hardy and Max Xu for discussing their works which are closely related to this problem. Some of this work was completed while the author was visiting the Centre de Recherches de Math\'{e}matiques de Montr\'{e}al and thanks them for excellent working conditions. The author is supported by the Heilbronn Institute for Mathematical Research.
\section{Preliminary Results} \label{sec2}
\subsection{A partition of unity}
Before we can get to introducing the most important prerequisites, we need a suitable partition of unity. This will allow us to transfer between estimates over character sums and estimates for sums of Steinhaus random multiplicative functions.
\begin{lemma}
    (Approximation Result 1 in \cite{harper2023typicalsizecharacterzeta}) Let $N$ be large and $\delta>0$ be small. Then there exists functions $g,g_{N+1}:\mathbb{R} \to \mathbb{R}$ (depending on $N$ and $\delta$) such that for $g_j(x) = g(x-j)$ where $j$ are all the integers such that $|j| \leq N$, the following hold:
    \begin{enumerate}
        \item $\sum_{|j| \leq N}g_j(x) + g_{N+1}(x) = 1, \quad \forall x \in \mathbb{R}$;
        \item $g(x) \geq 0, \quad \forall x \in \mathbb{R}$ and $g(x) \leq \delta$ whenever $|x| \geq 1$;
        \item $g_{N+1}(x) \geq 0, \quad \forall x \in \mathbb{R}$ and $g_{N+1}(x) \leq \delta$ whenever $|x| \leq N$;
        \item $\forall \ell \in \mathbb{N}$ and $\forall x \in \mathbb{R}$, the derivative estimate $\big|\frac{d^{\ell}}{dx^{\ell}}(g(x))\big| \leq \frac{1}{\pi(\ell+1)}\bigg(\frac{2\pi}{\delta}\bigg)^{\ell+1}$ holds.
    \end{enumerate}
    \label{partition_result}
\end{lemma}
\subsection{Estimates involving character averaging}
The most important tools in the proof are estimates involving character sums and connecting suitably averaged sums to their random multiplicative counterparts. These are proved in \cite{harper2023typicalsizecharacterzeta}. We start with an even moment estimate.
\begin{lemma}
    (Lemma 1 in \cite{harper2023typicalsizecharacterzeta}) Let $x \geq 1$ and $(c(n))$ be a set of arbitrary complex numbers. Let $\mathcal{P}$ be a finite set of primes, and $\mathcal{R}$ be a set containing the elements of $\mathcal{P}$ and their squares. Let $U = \max\{r \in \mathcal{R}\}$. Finally, set $R(\chi) = \sum_{r \in \mathcal{R}}\frac{a(r)\chi(r)}{\sqrt{r}}$ with $(a(r))$ a separate set of arbitrary complex numbers. Then for any $k \in \mathbb{N}$ such that $xU^k < q$
    $$\mathbb{E}^{\text{char}}\big|\sum_{n \leq x}c(n)\chi(n)\big|^2|R(\chi)|^{2k} \ll \big(\sum_{n \leq x}\tilde{d}(n)|c(n)|^2\big)\cdot(k!)\big(\sum_{r \in \mathcal{R}}\frac{2v_r|a(r)|^2}{r}\big)^k$$
    where $\tilde{d}(n) := \sum_{d|n}\mathbf{1}_{p|d \Rightarrow p \in \mathcal{P}}$ and $v_r$ being equal to $1$ if $r$ is a prime, and $6$ otherwise. \label{even_moment}
\end{lemma}
This will also act as a moment estimate when decoupling the log-normal and log-correlated regimes in the mesoscopic case. The key result is the following randomisation estimate which allows us to pass from character averages to averages over Steinhaus multiplicative functions. It will allow us to approximate our barrier events as smooth functions which is very beneficial in this work. This will act as a condensed version of much of the work done in Key Proposition 1 in \cite{harper2019partition}.
\begin{lemma}
     (Proposition 1 in \cite{harper2023typicalsizecharacterzeta}) Let $g_j$ with associated parameters $N,\delta$ be as described in Lemma \ref{partition_result}. Further, let $x \geq 1$ and suppose $(c(n))_{n \leq x}$ are arbitrary complex coefficients such that $|c(n)| \leq 1$ for each $n$. Finally, suppose $P$ is large and $Y \in \mathbb{N}$ such that $xP^{400(Y/\delta)^2\log(N\log(P))} < q$. \par
    Then for any indices $-N \leq j(1),...,j(Y)\leq N$ and sequences $((a_i(p),a_i(p^2))_{p \leq P})_{i \leq Y}$ of complex numbers with absolute value $\leq 1$, we have
    \begin{align*}
        \mathbb{E}^{\text{char}}\prod_{i=1}^Y&g_{j(i)}\big(\Re\big(\sum_{p \leq P}\frac{a_i(p)\chi(p)}{p^{1/2}} + \frac{a_i(p^2)\chi(p^2)}{p}\big)\big)\big|\sum_{n \leq x}\frac{c(n)\chi(n)}{n^{1/2}}\big|^2 \\
        = \mathbb{E}\prod_{i=1}^Y&g_{j(i)}\big(\Re\big(\sum_{p \leq P}\frac{a_i(p)f(p)}{p^{1/2}} + \frac{a_i(p^2)f(p^2)}{p}\big)\big)\big|\sum_{n \leq x}\frac{c(n)f(n)}{n^{1/2}}\big|^2 + O\big(\frac{\log(x)}{(N\log(P))^{Y/\delta^2}}\big).
    \end{align*} \label{randomisation}
\end{lemma}
\begin{proof}
This is identical to that of Proposition 1 in \cite{harper2023typicalsizecharacterzeta}, but we can estimate $\sum_{n \leq x}\frac{|c(n)|^2\tilde{d}(n)}{n}$ in a sharper manner. Indeed, we observe that by Lemma \ref{even_moment}
$$\sum_{n \leq x}\frac{|c(n)|^2\tilde{d}(n)}{n} \leq \sum_{\substack{m \leq x \\ P(m) \leq P}}\frac{d(m)}{m}\sum_{\substack{n \leq x/m \\ P(n) > P}}\frac{1}{n} \ll \log(x)\log(P).$$
Following the rest of the calculations in the proposition yields the result.
\end{proof}
Note that virtually all of the choices of $\delta$ and $P$ are motivated by the condition relating the various parameters above, so it is in our interest to make sure both $P$ and $\delta$ are as large as possible while getting this result to hold. Next we need some results on the Dirichlet $L$-functions on the critical line.
\begin{lemma}
    (Theorem 1 in \cite{davies1965}) Let $q$ be a large prime, and $\chi$ be a Dirichlet character modulo $q$. Then for any $|h| \leq 1$
    $$L(1/2+ih,\chi) = \sum_{n \leq \sqrt{q/2\pi}}\frac{\chi(n)}{n^{1/2+ih}} + \psi(1/2+it,\chi)\sum_{n \leq \sqrt{q/2\pi}}\frac{\overline{\chi(n)}}{n^{1/2-ih}} + O(q^{-1/4}),$$
    where $|\psi(1/2+it,\chi)|=1$ for any $\chi$ modulo $q$.
\label{Dirichlet_AFE}
\end{lemma}
Away from the central point, one then has the length of the sum we consider to be $qt$ (corresponding to $L(1/2+it,\chi)$) with the error term of $(qt)^{-1/4}$ (seen in the exercises in \cite{koukoulopoulos2019distribution}). Consequently our results could be recovered away from the central point as the parameters we choose for our results will allow for Lemma \ref{randomisation} to work. With this in hand, we want to introduce a suitable approximation when we average over all characters. We do this to shorten the Dirichlet polynomial in the above, as well as to connect more easily to our random sums. This is our first step to separate the log normal part of the Dirichlet polynomial and the log correlated part. This where we identify the hard boundary for mollifying $|L(1/2+ih,\chi)|$ directly in the mesoscopic case. No matter how we vary $\varepsilon_\theta$, we simply can't extract enough cancellation to prove meaningful results when $\theta$ is slightly larger than $-1/2$ (one can turn this into a quantitative result, but we choose against doing this for sake of the exposition).
\begin{lemma}
    Fix $\theta \in (-1/2,0]$. Then uniformly for all $|\sigma-1/2| \leq \frac{1}{\log(q)}$ and $|h| \leq \log^{\theta}(q)$
    \begin{equation}
        \mathbb{E}^{\text{char}}\big|\sum_{n \leq \sqrt{q/2\pi}}\frac{\chi(n)}{n^{\sigma+ih}}-\sum_{\substack{m \leq q^{\varepsilon_\theta} \\ P^+(m) \leq P_\theta}}\frac{\chi(m)}{m^{1/2+ih}}\sum_{\substack{n \leq q^{1/2-2\varepsilon_\theta} \\ P^-(n) > P_\theta}}\frac{\chi(n)}{n^{\sigma+ih}}\big|^2 \ll \frac{\log^{1+\theta}(q)}{(\log\log(q))^2}.  \label{factored_eqn}
    \end{equation} \label{factored_form}
\end{lemma}
\begin{proof}
    We use the triangle inequality to handle the lack of shift in the smooth part of the sum. This means the left hand side of equation (\ref{factored_eqn}) satisfies
    \begin{align*}
        \ll& \mathbb{E}^{\text{char}}\big|\sum_{n \leq \sqrt{q/2\pi}}\frac{\chi(n)}{n^{\sigma+ih}}-\sum_{\substack{m \leq q^{\varepsilon_\theta} \\ P^+(m) \leq P_\theta}}\frac{\chi(m)}{m^{\sigma+ih}}\sum_{\substack{n \leq q^{1/2-2\varepsilon_\theta} \\ P^-(n) > P_\theta}}\frac{\chi(n)}{n^{\sigma+ih}}\big|^2 \\
        &+ \mathbb{E}^{\text{char}}\big|\sum_{\substack{n \leq q^{1/2-2\varepsilon_\theta} \\ P^-(n) > P_\theta}}\frac{\chi(n)}{n^{\sigma+ih}}\big|^2\mathbb{E}^{\text{char}}\big|\sum_{\substack{m \leq q^{\varepsilon_\theta} \\ P^+(m) \leq P_\theta}}\frac{\chi(m)}{m^{\sigma+ih}}-\sum_{\substack{m \leq q^{\varepsilon_\theta} \\ P^+(m) \leq P_\theta}}\frac{\chi(m)}{m^{1/2+ih}}\big|^2.
    \end{align*}
    We achieve the second term by using the orthogonality of Dirichlet characters. For the first term, we are counting the numbers that are between $q^{1/2-2\varepsilon_\theta}$ and numbers with a $P$-smooth part which is greater than $q^{\varepsilon_\theta}$. Then by Lemma \ref{even_moment} (or directly by orthogonality), the first term is bounded by
    \begin{align*}
        &\sum_{q^{1/2-2\varepsilon_\theta} \leq n \leq q^{1/2}}  \frac{1}{n^{2\sigma}}+\sum_{\substack{q^{\varepsilon_\theta}< m \leq q^{1/2-2\varepsilon_\theta} \\ P^+(n) \leq P_\theta}}\frac{1}{m^{2\sigma}}\sum_{\substack{n \leq q^{1/2-2\varepsilon_\theta} \\ P^-(n) > P_\theta}}\frac{1}{n^{2\sigma}} \\
        \ll& \varepsilon_\theta\log(q) + \frac{\log(q)}{\log(P_\theta)}\sum_{\substack{ m > q^{\varepsilon_\theta} \\ P^+(n) \leq P_\theta}}\frac{1}{m^{2\sigma}} \\ \ll& \log(q)(\varepsilon_\theta+e^{-\frac{\varepsilon_\theta\log(q)}{\log(P_\theta)}}) \ll \frac{\log^{1+\theta}(q)}{(\log\log(q))^2} + \log(q)\exp\big(-\frac{\log^{1+2\theta}(q)}{(\log\log(q))^2}\big).
    \end{align*}
    The final line is achieved by using Rankin's trick to bound the contribution of the smooth numbers by $\log(P_\theta)e^{-\frac{\varepsilon_\theta\log(q)}{\log(P_\theta)}}$ and by noting $\theta > -1/2$. As for the second term, this reduces to bounding
    $$\frac{\log(q)}{\log(P_\theta)}\sum_{\substack{m \leq q^{\varepsilon_\theta} \\ P^+(m) \leq P_\theta}}\frac{|m^{\sigma-1/2}-1|^2}{m}.$$
    Observing that $|\sigma-1/2| \leq \frac{1}{\log(q)}$, then $|e^{\log(m)|\sigma-1/2|}-1| \leq \frac{\log(m)}{\log^{1+\theta}(q)}$ by Taylor expansion, so the contribution of the second term is $\ll \varepsilon_\theta^2\log^{1+\theta}(q)$.
\end{proof}
To handle the contribution of log-correlated part of the Dirichlet polynomial, we will need to further mollify $\sum_{\substack{n \leq q^{1/2-2\varepsilon_\theta} \\ P^{-}(n) > P_\theta}}\frac{\chi(n)}{q^{1/2+ih}}$ in a suitable sense. This is just an altered version of Lemma \ref{factored_form}, and consequently its proof is omitted as it is an application of orthogonality. This lemma holds in its stated form, or when the $\chi$ are replaced by Steinhaus random multiplicative functions $f(n)$.
\begin{lemma}
    Fix $\theta \in (-1/2,0]$. Suppose $P > P_\theta$ and $\min\{\frac{1}{\log(q)},\varepsilon_\theta\} < \varepsilon_1 \leq \frac{1}{10}$ say, then we have uniformly for $|\sigma-1/2| \leq \frac{1}{\log(q)}$ and $|h| \leq \log^{\theta}(q)/2$
    \begin{equation*}
    \mathbb{E}^{\text{char}}\big|\sum_{\substack{n \leq q^{1/2-2\varepsilon_\theta} \\ P^{-}(n) >P_\theta}}\frac{\chi(n)}{n^{\sigma+ih}}-\sum_{\substack{m \leq q^{\varepsilon} \\ p|m \Rightarrow P_\theta < p \leq P}}\frac{\chi(m)}{m^{1/2+ih}}\sum_{\substack{n \leq q^{1/2-2\varepsilon} \\ P^{-}(n) > P}}\frac{\chi(n)}{n^{\sigma+ih}}\big|^2 \ll \log^{1+\theta}(q)\big(\varepsilon+e^{-\frac{\varepsilon\log(q)}{\log(P)}}\big),
    \end{equation*}
    where $\varepsilon = \varepsilon_\theta+\varepsilon_1$. \label{logcor_factored_form}
\end{lemma}
The next result is needed only for the proof of Theorem \ref{short_range_max}, which is on a conditional fourth moment estimate. This effectively comes down to saying that $L(1/2+ih,\chi)$ is large, then $\sum_{\substack{m \leq q^\varepsilon \\ P^+(m) \leq P}}\frac{\chi(m)}{m^{1/2+ih}}$ must be large too. We let
$$\mathcal{H}_\chi(h;V) := \big\{\big|\sum_{\substack{m \leq q^\varepsilon \\ P^+(m) \leq P}}\frac{\chi(m)}{m^{1/2+ih}}\big| \leq \frac{\log(P)}{V}\big\}.$$
\begin{lemma}
    Uniformly for large $ P \leq \sqrt{q}, \ 0 <\varepsilon < 1/10$ with $V>0$ and $|\sigma-1/2| \leq 1/\log(q)$, we have
    $$\mathbb{E}^{\text{char}}\big(\mathbf{1}_{\mathcal{H}_\chi(h;V)}(h)\big|\sum_{\substack{m \leq q^\varepsilon \\ P^+(m) \leq P}}\frac{\chi(m)}{m^{1/2+ih}}\sum_{\substack{n \leq q^{1/2-2\varepsilon} \\ P^{-}(n) > P}}\frac{\chi(n)}{n^{\sigma+it}}\big|^4\big) \ll \frac{1}{V^2}\frac{\log^4(q)}{\log(P)}.$$ \label{conditional_fourth}
\end{lemma}
\begin{proof}
    Using the conditioning given by the event $\mathcal{H}_\chi(h;V)$, we can upper bound the left hand side by
    \begin{align*}
        &\big(\frac{\log(P)}{V}\big)^2\mathbb{E}^{\text{char}}\big(\big|\sum_{\substack{m \leq q^\varepsilon \\ P^+(m) \leq P}}\frac{\chi(m)}{m^{1/2+ih}}\big|^2\big|\sum_{\substack{n \leq q^{1/2-2\varepsilon} \\ P^{-}(n) > P}}\frac{\chi(n)}{n^{\sigma+it}}\big|^4\big) \\
        \ll& \big(\frac{\log(P)}{V}\big)^2\mathbb{E}^{\text{char}}\big(\big|\sum_{\substack{m \leq q^\varepsilon \\ P^+(k) \leq P}}\frac{\chi(m)}{m^{1/2+ih}}\big|^2\big)\big|\sum_{\substack{n \leq q^{1-4\varepsilon} \\ P^{-}(n) > P}}\frac{c(n)}{n^{2\sigma}}\big|^2,
    \end{align*}
    where $c(n) := \#\{(d_1,d_2): d_1,d_2 \leq q^{1/2-2\varepsilon},  n=d_1d_2\}$ (where we have applied Lemma \ref{even_moment} to the rough sum due to orthogonality of the characters). Then Lemma \ref{even_moment} implies that this
    $$\ll \big(\frac{\log(P)}{V}\big)^2\sum_{P^+(m) \leq P}\frac{1}{m}\sum_{\substack{n \leq q^{1-4\varepsilon} \\ P^-(n) > P}}\frac{(c(n))^2}{n} \ll \frac{1}{V^2}\log^3(P)\sum_{\substack{n \leq q^{1-4\varepsilon} \\ P^-(n) > P}}\frac{d_4(n)}{n}$$
    where $d_k(n) = \#\{(d_1,...,d_k):\prod_{i=1}^kd_i=n\}$ the standard $k$-fold divisor function (we used that $(d_2(n))^2 \leq d_4(n)$ here). The final term can be bounded \\ by $\prod_{P < p \leq q^{1-4\varepsilon}}\big(1-\frac{4}{p}\big)^{-1} \ll \big(\frac{\log(q)}{\log(P)}\big)^4$, which yields the claim.  
\end{proof}
\subsection{Probabilistic results} \label{prob_heur}
While we will not explicitly use these results in this paper (since we state the adaptations we need to make to Lemmas 4 and 7 in \cite{harper2019partition} and treat them as blackboxes), we will still state the two probabilistic results which are the heavy machinery behind much of the work going on in this paper. The first of which is behind a lot of the proof of Theorems \ref{pf_moments_constant} and \ref{pf_moments_meso}.
\begin{lemma}
    (Probability Result 1, \cite{aharper201}) Let $a\geqslant 1$. For any sufficiently large integer $n > 1$, let $(G_k)_{k=1}^n$ be a sequence of independent real Gaussian random variables with mean $0$ and variance between $1/20$ and $20$, say. Suppose $h$ is a function such that $|h(j)| < 10\log(j)$. Then
    $$\Prob(\sum_{m=1}^j G_m \leqslant a + h(j),  \ \forall \ 1 \leqslant j \leqslant n ) \asymp \min\big\{\frac{a}{\sqrt{n}},1\big\}.$$ \label{ballot}
\end{lemma}
The second is a more specialised result which is designed for finding unusually large values.
\begin{lemma}
    (Probability Result 1, \cite{harper2019partition}) Let $a,b$ and $n \in \mathbb{N}$ be sufficiently large, with $(G_k)_{k \leq n}$ be a sequence of independent Gaussian random variables with mean $0$ and variance between $1/20$ and $20$ (say). Then we have the uniform upper bound
    $$\mathbb{P}(\sum_{k=1}^jG_k \leq a \ \forall1 \leq j \leq n, \text{ and } a-b \leq \sum_{k=1}^nG_k \leq a) \ll \min\big\{1,\frac{a}{\sqrt{n}}\big\}(\min\big\{1,\frac{b}{\sqrt{n}}\big\})^2.$$
\end{lemma}
To adapt these into the random Euler product setting, Harper derives Lemmas 4 and 5 in \cite{aharper201} to do this. Since their statements are rather long, we choose against doing this, but will try to explain the differences between the pure probability situation and the random Euler product one. To do this, we set for $0 \leq \ell \leq \log\log(P)$ for some large $P$
$$I_{\ell,\text{rand}}(h):=\prod_{P^{e^{-(\ell+1)}}< p \leq P^{e^{-\ell}}}\big(1-\frac{f(p)}{p^{1/2+ih}}\big)^{-1}.$$
It turns out (due a direct Gaussian comparison) that $\log(|I_{\ell,\text{rand}}(h)|)$ is well approximated by a standard Gaussian random variable of mean $1$ and variance $1$. This means that $\log(|F_P(1/2+ih)|) := \sum_{\ell \leq \log\log(P)}\log(|I_{\ell,\text{rand}}(h)|)$ can be seen as a sum of Gaussians. This strategy works quite well, except for very large $\ell$ values which correspond to very small primes (which are not very Gaussian). Consequently, we introduce a large constant $B$ and we can show that $\log(|F_p(s)|) \approx \sum_{\ell \leq \log\log(P)-B-1}\log(|I_{\ell,\text{rand}}(h)|)$, and proceed with applying pure probability results like the ones stated above to this situation. This is a very high level interpretation of what is really going on, so if we need to use these results then we will recall them from \cite{aharper201}.  
\section{The chaining argument for lognormal contribution} \label{chain_arg}
In this section we will prove a chaining argument which we will use throughout the paper. The log-normal contribution is such a short Dirichlet polynomial, we can afford to approximate by a truncated Euler product once we have transferred from character sums to random multiplicative functions. 
\begin{lemma}
    Let $\theta \in (-1/2,0]$. Then uniformly for all $|h| \leq \log^{\theta}(q)/2$
    \begin{equation}
        \mathbb{E}\big|\sum_{\substack{m \leq q^{\varepsilon_\theta} \\ P^+(m) \leq P_\theta}}\frac{f(m)}{m^{1/2+ih}} - \exp\big(\sum_{p^j \leq P_\theta}\frac{f(p)^j}{jp^{j(1/2+ih)}}\big)\big|^2\ll \log(P_\theta)e^{-\frac{\varepsilon_\theta \log(q)}{\log(P_\theta)}} \label{exp_log_eq}
    \end{equation} 
    \label{exp_logn}
\end{lemma}
\begin{proof}
    We see that
    \begin{equation*}
        \begin{split}\mathbb{E}\big|\sum_{\substack{m \leq q^{\varepsilon_\theta} \\ P^+(m) \leq P_\theta}}\frac{f(m)}{m^{1/2+ih}} - \exp\big(\sum_{p^j \leq P_\theta}\frac{f(p)^j}{jp^{j(1/2+ih)}}\big)\big|^2 \ll \mathbb{E}\big|\sum_{\substack{m \leq q^{\varepsilon_\theta} \\ P^+(m) \leq P_\theta}}\frac{f(m)}{m^{1/2+ih}}-\sum_{\substack{P^+(m) \leq P_\theta}}\frac{f(m)}{m^{1/2+ih}}\big|^2 \\ +\mathbb{E}\big|\sum_{\substack{P^+(m) \leq P_\theta}}\frac{f(m)}{m^{1/2+ih}}-\exp\big(\sum_{p^j \leq P_\theta}\frac{f(p)^j}{jp^{j(1/2+ih)}}\big)\big|^2\end{split}
    \end{equation*}
    where we have used $|A+B|^2 \leq 2(|A|^2+|B|^2)$. The first term on the right can be then bounded by $\sum_{\substack{m > q^{\varepsilon_\theta} \\ P^+(m) \leq P_\theta}}\frac{f(m)}{m^{1/2+ih}}$, where we apply Rankin's trick in a similar fashion to that seen in Proposition \ref{factored_form} and we obtain our result. As for the second term on the right, notice that
    $$\sum_{\substack{P^+(m) \leq P_\theta}}\frac{f(m)}{m^{1/2+ih}} = \exp\bigg(\sum_{p^j :p \leq P_\theta}\frac{f(p)^j}{jp^{j(1/2+ih)}}\bigg),$$
    so we can rewrite the second term as
    $$\mathbb{E}\big|\exp\big(\sum_{p^j :p \leq P_\theta}\frac{f(p)^j}{jp^{j(1/2+ih)}}\big)-\exp\big(\sum_{p^j \leq P_\theta}\frac{f(p)^j}{jp^{j(1/2+ih)}}\big)\big|^2.$$
    After splitting the sum, we have this is
    $$\mathbb{E}\big|\exp\big(\sum_{p^j \leq P_\theta}\frac{f(p)^j}{jp^{j(1/2+ih)}}\big)\big(\exp\big(\sum_{p^j >P_\theta:p \leq P_\theta}\frac{f(p)^j}{jp^{j(1/2+ih)}}\big)-1\big)\big|^2.$$
    The remaining terms in the second exponential contribute at most $O(\frac{1}{\sqrt{P_\theta}\log(P_\theta)})$, so we Taylor expand the second exponential to obtain this is
    $$\ll \frac{1}{\sqrt{P_\theta}\log(P_\theta)}\mathbb{E}\big|\exp\big(\sum_{p^j \leq P_\theta}\frac{f(p)^j}{jp^{j(1/2+ih)}}\big)\big|^2 \ll \frac{1}{\sqrt{P_\theta}}$$
    which is certainly small enough.
\end{proof}
For simplicity, we discard all the $j \geq 3$ as all of the terms will then be summable for the rest of this section. As for the chaining argument, it follows from the next two lemmas. We define the events
\begin{equation}
    \mathcal{S}_\chi(A) :=\big\{\max_{|h| \leq \log^{\theta}(q)}\big|\sum_{p \leq P_\theta}\bigg(\sum_{v=1,2}\frac{\chi(p^v)}{vp^{v(1/2+ih)}}-\frac{\chi(p^v)}{p^{v}}\bigg) \big| \leq A\big\}, \label{central_char_shift}
\end{equation}
and its random counterpart
\begin{equation}
    \mathcal{S}_{\text{rand}}(A) :=\big\{\max_{|h| \leq \log^{\theta}(q)}\big|\sum_{p \leq P_\theta}\sum_{v=1,2}\big(\frac{f(p^v)}{vp^{v(1/2+ih)}}-\frac{f(p^v)}{p^{v}}\big) \big| \leq A\big\}. \label{central_rand_shift}
\end{equation}
Before we proceed to studying $\mathcal{S}_\chi(A)$ in more detail, we need to discretise the interval $|h| \leq \log^{\theta}(q)/2$. This will reduce our tasks involving $\mathcal{S}_\chi(A)$ to investigating the event only at a discrete set of points. For the next lemma, we introduce the notation
$$Y_{x,\chi}(h) = \sum_{p \leq x}\big(\frac{\chi(p)}{p^{1/2+ih}}+\frac{(\chi(p))^2}{p^{1+2ih}}\big).$$ Then the difference $Y_{P_\theta,\chi}(h)-Y_{P_\theta,\chi}(0)$ is also a Dirichlet polynomial, so we can apply the following lemma from \cite{abr20} (which is also Lemma A.1 in \cite{ah24}).
\begin{lemma}
    (Lemma 27 in \cite{abr20}) Let $D(s,\chi) = \sum_{n \leq N}\frac{\chi(n)}{n^s}$, a Dirichlet polynomial of length $N$. Then for any $\gamma \geq 100$
    \begin{equation*}\begin{split}\max_{|h| \leq \log^\theta(q)/2}|D(1/2+ih,\chi)|^2 \ll&\sum_{|j| \leq 16\log^\theta(q)\log(N)}\big|D\big(\frac{1}{2}+\frac{2\pi i j}{8\log(N)},\chi \big)\big|^2 \\
    &+ \sum_{|j| > 16\log^\theta(q)\log(N)}\frac{1}{1+|j|^\gamma}\big|D\big(\frac{1}{2}+\frac{2\pi i j}{8\log(N)},\chi \big)\big|^2.\end{split}
    \end{equation*}
\end{lemma}
Since Lemma 27 in \cite{abr20} is actually stated for the maximum over intervals of length $2$, this goes through without issue (since it holds for a general Dirichlet polynomial). Finally, we will show that $\mathcal{S}_\chi(A)$ occurs with high probability. This follows in an identical fashion to results in \cite{ah24} (only we do this in the discrete setting).
\begin{lemma}
    Let $\theta \in (-1/2,0]$. Then there exists an absolute $c>0$ such that the following is true. Suppose $0 \leq A\leq \log^{(1+\theta)/2}(q)$, then
    $$\mathbb{P}_q\big(\mathcal{S}_\chi(A) \text{ fails}\big) \ll e^{-A^2/c}.$$ \label{char_chaining}
\end{lemma}
\begin{proof}
    We begin by introducing a discretisation over the positions of $h$. We see that for $h_j = \frac{\pi j \log^{\theta}(q)}{4}$, then
    \begin{align*}
        &\mathbb{P}_q\big(\max_{|h| \leq \log^{\theta}(q)/2}\big|Y_{P_\theta,\chi}(h)-Y_{P_\theta,\chi}(0) \big| > A\big) \numberthis \label{discrete_lognormal} \\
        \ll& \mathbb{P}_q\big(\sum_{|j| \leq 32}|Y_{P_\theta,\chi}(h_j)-Y_{P_\theta,\chi}(0)|^2 > A^2/C\big) + \mathbb{P}_q\big(\sum_{|j| > 32}\frac{|Y_{P_\theta,\chi}(h_j)-Y_{P_\theta,\chi}(0)|^2}{1+|j|^{100}} > A^2/C\big)
    \end{align*}
    for $C>1$ an arbitrary fixed constant. This follows from choosing $\gamma=100$ and $N=P_\theta^2$ in the previous lemma. To handle the first sum, we notice
    \begin{align*}
        &\mathbb{P}_q\big(\sum_{|j| \leq 32}|Y_{P_\theta,\chi}(h_j)-Y_{P_\theta,\chi}(0)|^2 > A^2/C\big) \ll \sum_{|j| \leq 32}\mathbb{P}_q\big(|Y_{P_\theta,\chi}(h_j)-Y_{P_\theta,\chi}(0)|^2 > A^2/C\big),
    \end{align*}
    and apply Markov's inequality to the $r=\lfloor \frac{A^2}{2048C\pi^2}\rfloor$ (which is permitted by Lemma \ref{even_moment} as $P_\theta^r < q$ as $P_\theta^r \leq q^{1/1000}$ say) to find for each individual $j$
    \begin{align*}
        \mathbb{E}^{\text{char}}\bigg(|Y_{P_\theta,\chi}(h_j)-Y_{P_\theta,\chi}(0)|^2 > A^2/C\bigg)
        \ll& \frac{\mathbb{E}^{\text{char}}(|Y_{P_\theta,\chi}(h_j)-Y_{P_\theta,\chi}(0)|^{2r})}{(A^2/C)^{r}} \\ \ll& \frac{(r!)\big(\sum_{p \leq P_\theta}\frac{2|p^{-ih_j}-1|^2}{p}+\frac{12|p^{-2ih_j}-1|^2}{p^2}\big)^r}{(A^2/C)^r} \\
        \ll& \frac{2^r(r!)\big(\sum_{p \leq P_\theta}\frac{|p^{-ih_j}-1|^2}{p}+O(1)\big)^r}{(A^2/C)^r} \\
        \ll& \frac{2^r(r!)(\pi j/4)^{r}}{(A^2/C)^{r}}.
    \end{align*}
    Applying the Stirling approximation for $(r!)$ (that for $r$ sufficiently large, $r! \sim \sqrt{2\pi r}(r/e)^r$) then shows
    \begin{equation}
        \sum_{|j| \leq 32}\frac{2^r(r!)(\pi j/4)^{2r}}{(A^2/C)^{r}} \ll e^{-A^2/c} \label{small_j}
    \end{equation}
    for some constant $c>0$. To handle the second term in equation \ref{discrete_lognormal}, we observe
    \begin{align*}
        \big\{ \sum_{|j| > 32}\frac{|Y_{P_\theta,\chi}(h_j)-Y_{P_\theta,\chi}(0)|^2}{1+|j|^{100}} > A^2/C\big\} \subseteq \bigcup_{|j| > 32}\big\{|Y_{P_\theta,\chi}(h_j)-Y_{P_\theta,\chi}(0)|^2 > \frac{A^2|j|^{50}}{C}\big\}.
    \end{align*}
    From this, similar to before, we have
    \begin{align*}
        &\mathbb{P}_q\big(\sum_{|j| > 32}|Y_{P_\theta,\chi}(h_j)-Y_{P_\theta,\chi}(0)|^2 > A^2/C\big) \ll \sum_{|j| > 32}\mathbb{P}_q\big(|Y_{P_\theta,\chi}(h_j)-Y_{P_\theta,\chi}(0)|^2 > \frac{A^2|j|^{50}}{C}\big).
    \end{align*}
    Then we apply Markov's inequality to the $r=\lfloor\frac{32A^2}{\pi^2C}\rfloor$, so for each $j$ we have
    \begin{align*}
        \mathbb{P}_q\big(|Y_{P_\theta,\chi}(h_j)-Y_{P_\theta,\chi}(0)|^2 > \frac{A^2|j|^{50}}{C}\big) \ll& \frac{\mathbb{E}^{\text{char}}\big|Y_{P_\theta,\chi}(h_j)-Y_{P_\theta,\chi}(0)\big|^{2r}}{(A^2|j|^{50}/C)^r} \\
        \ll& \frac{\mathbb{E}^{\text{char}}(|Y_{P_\theta,\chi}(h_j)-Y_{P_\theta,\chi}(0)|^{2r})}{(A^2||j|^{50}/C)^{r}} \\ \ll& \frac{(r!)\big(\sum_{p \leq P_\theta}\frac{2|p^{-ih_j}-1|^2}{p}+\frac{12|p^{-2ih_j}-1|^2}{p^2}\big)^r}{(A^2|j|^{50}/C)^r} \\
        \ll& \frac{2^r(r!)(\pi/4)^{2r}}{(A^2|j|^{40}/C)^{r}}
    \end{align*}
    using Lemma \ref{even_moment}. Summing over the $j$ followed by the Stirling approximation for $(r!)$ with equation (\ref{small_j}) is sufficient to conclude
    $$\mathbb{P}_q\big(\max_{|h| \leq \log^{\theta}(q)}\big|Y_{P_\theta,\chi}(h)-Y_{P_\theta,\chi}(0) \big| > A\big) \ll e^{-A^2/c}.$$
\end{proof}
We can finally demonstrate that this is sufficient for our chaining bound. While we never apply this proposition in this form, we will use the part of the proof that involves random multiplicative functions many times and this is one of the cleaner ways to state the result. In either case we will refer to the below proposition as and when we require it in either the random or deterministic setting.
\begin{proposition}
    (Single point conditioning) Fix $\theta \in (-1/2,0]$. Then uniformly
    \begin{align*}
        \int_{|h| \leq \log^\theta(q)/2}\mathbb{E}_W^{\text{char}}|\sum_{\substack{m \leq q^{\varepsilon_\theta} \\ P^+(m) \leq P_\theta}}\frac{\chi(m)}{m^{1/2+ih}}|^2dh \ll e^{2W}\log^{\theta}(q)
    \end{align*} \label{chaining}
\end{proposition}
    \begin{proof}
    We introduce the event $\mathcal{T}_\chi(A)$ for a suitably large constant $A>0$ where
    $$\mathcal{T}_\chi := \mathcal{S}_{\chi}(A) \cap \{Y_{P_\theta,\chi}(0) \in [W,W+1]\}.$$
    Clearly, $\mathbf{1}_{\mathcal{T}_\chi(A)} \leq \mathbf{1}_{\mathcal{T}_\chi(h)}(h)$, which $\mathcal{T}_\chi(h)$ denotes that the event $\mathcal{S}_\chi(A)$ occurs for a particular $h$ (and not all $h$ simultaneously) and $Y_{P_\theta,\chi}$ is in a given interval $[W,W+1]$. Then since our integral is over a finite length, we can approximate $\mathbf{1}_{\mathcal{T}_\chi(h)}(h)$ by a smooth function $g(h)$ which satisfies various properties stated in Lemma \ref{partition_result}. At this point, we have shown that for some $\delta_s,N_s$ to be chosen later that
    \begin{align*}
        &\int_{|h| \leq \log^{\theta}(q)/2}\mathbb{E}^{\text{char}}\big(\mathbf{1}_{\mathcal{T}_\chi(A)}|\sum_{\substack{m \leq q^{\varepsilon_\theta} \\ P^+(m) \leq P_\theta}}\frac{\chi(m)}{m^{1/2+ih}}|^2\big)dh \\ \leq& \int_{|h| \leq \log^{\theta}(q)/2}\mathbb{E}^{\text{char}}\big(g\big(\sum_{p \leq P_\theta}\sum_{v=1,2}\frac{\chi(p^v)}{vp^{v(1/2+ih)}}\big)|\sum_{\substack{m \leq q^{\varepsilon_\theta} \\ P^+(m) \leq P_\theta}}\frac{\chi(m)}{m^{1/2+ih}}|^2\big)dh.
    \end{align*}
    At this point, we can apply Lemma \ref{randomisation} to allow us to switch to a statement involving random multiplicative functions. This is very important because it allows us to compare this Dirichlet polynomial to a true truncated Euler product as there are no periodicity conditions to worry about. As a consequence, we have
    \begin{align*}
        &\int_{|h| \leq \log^{\theta}(q)/2}\mathbb{E}^{\text{char}}\big(g\big(\sum_{p \leq P_\theta}\sum_{v=1,2}\frac{\chi(p^v)}{vp^{v(1/2+ih)}}\big)|\sum_{\substack{m \leq q^{\varepsilon_\theta} \\ P^+(m) \leq P_\theta}}\frac{\chi(m)}{m^{1/2+ih}}|^2\big)dh \\=&\int_{|h| \leq \log^{\theta}(q)/2}\mathbb{E}\big(g\big(\sum_{p \leq P_\theta}\sum_{v=1,2}\frac{f(p^v)}{vp^{v(1/2+ih)}}\big)|\sum_{\substack{m \leq q^{\varepsilon_\theta} \\ P^+(m) \leq P_\theta}}\frac{f(m)}{m^{1/2+ih}}|^2\big)dh + O\big(\frac{\log(x)}{(N_s\log(P_\theta))^{1/\delta_s^2}}\big).
    \end{align*}
    Now we can undo the smoothing, noting that $0 \leq (g - \mathbf{1}_{\mathcal{T}_\chi(h)}(h)) \leq \delta_s$. Consequently, we have
    \begin{align*}
        &\int_{|h| \leq \log^{\theta}(q)/2}\mathbb{E}\big(g(\sum_{p \leq P_\theta}\sum_{v=1,2}\frac{f(p^v)}{vp^{v(1/2+ih)}})|\sum_{\substack{m \leq q^{\varepsilon_\theta} \\ P^+(m) \leq P_\theta}}\frac{f(m)}{m^{1/2+ih}}|^2\big)dh \\
        \leq& \int_{|h| \leq \log^{\theta}(q)/2}\mathbb{E}\big((\mathbf{1}_{\mathcal{T}_{\text{rand}}(h)}(h)+\delta_s)|\sum_{\substack{m \leq q^{\varepsilon_\theta} \\ P^+(m) \leq P_\theta}}\frac{f(m)}{m^{1/2+ih}}|^2\big)dh.
    \end{align*}
    At this point, we can apply the triangle inequality, and obtain the following upper bound on the above
    \begin{align*}
        \ll& \int_{|h| \leq \log^\theta(q)/2}\mathbb{E}\big((1+\delta_s)\big|\sum_{\substack{m \leq q^{\varepsilon_\theta} \\ P^+(m) \leq P_\theta}}\frac{f(m)}{m^{1/2+ih}}-\exp\bigg(\sum_{j=1,2}\sum_{p \leq P_\theta}\frac{f(p)^j}{jp^{j(1/2+ih)}}\bigg)\big|^2\big)dh \\
        &+\int_{|h| \leq \log^\theta(q)/2}\mathbb{E}\big((\mathbf{1}_{\mathcal{T}_{\text{rand}}(h)}(h)+\delta_s)\exp\big(2\sum_{j=1,2}\sum_{p \leq P_\theta}\frac{\Re((f(p)p^{-ih})^j)}{jp^{j/2}}\big)\big)dh.
    \end{align*}
    From Lemma \ref{exp_logn}, the first term's contribution is negligible as $\mathbf{1}_{\mathcal{T}_{\text{rand}}(h)}(h) \leq 1$. This leaves us to handle the term involving the exponential. We have $$\delta_s\int_{|h| \leq \log^\theta(q)/2}\mathbb{E}\big(\exp\big(2\sum_{j=1,2}\sum_{p \leq P_\theta}\frac{\Re((f(p)p^{-ih})^j)}{jp^{j/2}}\big)\big)dh\ll \delta_s,$$ using results on moments of random Euler products. The term involving $\mathbf{1}_{\mathcal{T}_{\text{rand}}(h)}(h)$ can then be further upper bounded by
    \begin{equation}
        \int_{|h| \leq \log^\theta(q)/2}\mathbb{E}\big(\mathbf{1}_{\mathcal{T}_{\text{rand}}(h)}(h)\exp\big(2\big(\sum_{j=1,2}\sum_{p \leq P_\theta}\frac{\Re((f(p)p^{-ih})^j)}{jp^{j(1/2)}}-Y_{P_\theta}(0)\big) + 2Y_{P_\theta}(0)\big)\big)dh. \label{final_form}
    \end{equation}
    We can then apply Cauchy Schwarz to the expectation, allowing us to split $\mathcal{T}_{\text{rand}}(h)$ back into $\mathcal{S}_{\text{rand}}(h)$ and conditioning on $Y_{P_\theta}(0)$ as $\mathbf{1}_{\mathcal{T}_{\text{rand}(h)}} \leq \mathbf{1}_{\mathcal{S}_{\text{rand}}(h)}\mathbf{1}_{Y_{P_\theta}(0) \in [W,W+1]}$. This shows the above satisfies
    \begin{equation*}\begin{split}
        &\int_{|h| \leq \log^\theta(q)/2}\bigg[\big(\mathbb{E}\big(\mathbf{1}_{\mathcal{S}_{\text{rand}}(h)}(h)\exp\big(4\big|\sum_{j=1,2}\sum_{p \leq P_\theta}\frac{\Re((f(p)p^{-ih})^j)}{jp^{j(1/2)}}-Y_{P_\theta}(0)\big|\big)\big)^{1/2} \\ \cdot&\big(\mathbb{E}_W[\exp(4Y_{P_\theta}(0))]\big)^{1/2}\bigg]dh.
        \end{split}
    \end{equation*}
    This is clearly bounded by
    $$\ll \log^{\theta}(q)e^{2(A+W)}.$$
    Now, to handle when $\mathcal{T}_\chi(A)$ does not occur, we apply a chaining argument, since we know that $$\mathbb{P}(\mathcal{T}_\chi(A)) \text{ fails}) \leq \sum_{k=1}^{4\log\log\log(P_\theta)}\mathbb{P}(\mathcal{T}_\chi(2^kA) \text{ but }\mathcal{T}_\chi(2^{k-1}A) \text{ fails})+\mathbb{P}(\mathcal{T}_\chi(((\log\log(P_\theta))^{4\log(2)}A) \text{ fails}).$$
    We can apply the same argument throughout to handle the terms in the sum, only this time when we arrive at the end, we will obtain the upper bound
    $$\ll \log^{\theta}(q)e^{-\frac{2^{2k}A^2}{4c} + 2^{k}A+2W}.$$
    Here we observe that $A(\log\log(P_\theta))^{4\log(2)} \ll A(\log\log(q))^{3}$, which for our choice of $A$, this is clearly $\leq \log^{(1+\theta)/2}(q)$ so Lemma \ref{char_chaining} holds. This ensures the sum is then $\ll \log^\theta(q)e^{2W}$. For the final term, then we can instead apply standard results for moments of random Euler products, which bounds the second term by
    $$\ll Ce^{-A^2(\log\log(P_\theta))^{2\log(2)}/c} =o(1)$$
    using that this decays faster than any power of $\log(q)$ as $q$ is large. We choose $A > \max\{1,\sqrt{c}\}$ a fixed constant, for simplicity.
    \end{proof}
\section{Proof of Theorems \ref{pf_moments_constant} and \ref{pf_moments_meso}} \label{pfm_proof}
This proof will act like a combination of \cite{ah24} and \cite{harper2019partition} for the different parts of the Dirichlet polynomial respectively. Since we have already proved Proposition \ref{chaining}, we can largely follow Harper's treatment of the problem for the log-correlated regime. Once we have handled this contribution, we then are free to extract the log-normal contribution. \par
For this purpose, suppose $P$ is a large quantity in terms of $q$ (to be chosen later) and fix $\theta \in (-1/2,0]$. We define the set of points for each $0 \leq j \leq \log\log(P)+\theta\log\log(q)-1$
\begin{equation}
    h(j) := \max\big\{u \leq h(j-1):u=\frac{n}{(\log(P)/e^j)\log(\log(P)/e^j)} \text{ for some } n \in \mathbb{Z}\big\} \label{discrete_h}
\end{equation}
with $h(-1) = h$. We then define the partial Euler products
\begin{equation}
    I_{\ell,\chi}(h) := \prod_{P^{e^{-(\ell+1)}} < p \leq P^{e^{-\ell}}}\big(1-\frac{\chi(p)}{p^{1/2+ih}}\big)^{-1} \label{char_increment}
\end{equation}
and its random equivalent
\begin{equation}
    I_{\ell,\text{rand}}(h) :=\prod_{P^{e^{-(\ell+1)}} < p \leq P^{e^{-\ell}}}\big(1-\frac{f(p)}{p^{1/2+ih}}\big)^{-1} \label{random_increment}
\end{equation}
where the $f(p)$ are independent Steinhaus random variables. Next, we define the functions $P_q(\theta) = \log(P)\log^\theta(q)$ and $g(j) = C\min\{\sqrt{\log(P_q(\theta))},\frac{1}{1-r}\} + 2\log\log(P_q(\theta)/e^j)$ where $C>0$ is a large absolute constant. Finally, we set $\mathcal{G}_\chi$ as the event such that for all $|h| \leq \log^\theta(q)/2$ and all $0 \leq j \leq \log(P_q(\theta)) - B -1$, we have
$$\big|\sum_{\ell=j}^{\lfloor\log(P_q(\theta))\rfloor-B-1}\log(|I_{\ell,\chi}(h(\ell))|)\big|\leq \log(P_q(\theta)) - j +g(j).$$
Note this is defined in exactly the same way as the event $\mathcal{G}_t$ in \cite{harper2019partition}, only we have $\chi$ in place of $t$ (with $B$ being a parameter required due to the discussion in Section \ref{prob_heur}). We define $\mathcal{G}_{\text{rand}}$ in almost the same way, with $f(p)$ a sequence of independent Steinhaus random variables being used in place of Dirichlet character $\chi$ modulo $q$, only where we use the function $g_{\text{rand}}(j) = g(j) + 19/3$ (the $19/3$ is used to absorb issues arising from extra convergent terms). We then have the following propositions, from which we may deduce Theorems \ref{pf_moments_constant} and \ref{pf_moments_meso}.
\begin{proposition}
    Fix $\theta \in (-1/2,0]$. Then uniformly for $\exp(\log^{1+\theta}(q)) \leq P \leq q^{1/(\log\log(q))^2}$ sufficiently large in terms of $C$ and $2/3 \leq r \leq 1$
    $$\mathbb{P}_q\big(\mathcal{G}_\chi \text{ fails}\big) \ll e^{-2C\min\{\sqrt{\log(P_q(\theta))},\frac{1}{1-r}\}}.$$ \label{bad_chaining_moment}
\end{proposition}
\begin{proposition}
    Fix $\theta \in (-1/2,0]$. Then uniformly for $\exp(\log^{|\theta|}(q)) \leq P \leq q^{1/(\log\log(q))^6}$ sufficiently large in terms of $C$ and \\ for $\max\{\frac{\log(P)\log\log\log(P)}{\log(q)},\varepsilon_\theta\}<\varepsilon < 1/10$, then uniformly for $|h| \leq \log^{\theta}(q)/2$ and $2/3 \leq r \leq 1$, we have
    \begin{align*}
        \mathbb{E}^{\text{char}}&\big(\mathbf{1}_{\mathcal{G}_\chi}|\sum_{\substack{m \leq q^\varepsilon \\ p|m \Rightarrow P_\theta < p \leq P}}\frac{\chi(m)}{m^{1/2+ih}}|^2|\sum_{\substack{n \leq q^{1/2-2\varepsilon} \\ P^{-}(n) > P}}\frac{\chi(n)}{n^{1/2+ih}}|^2dh\big)^r \\ \ll& \big(C\log^{1+\theta}(q)\min\{1,\frac{1}{(1-r)\sqrt{\log\log(P)}}\}\big)^r.
    \end{align*}\label{good_moment_bound}
\end{proposition}
We note that we will not use Proposition \ref{good_moment_bound} in this exact form, but it is close enough to the version we will apply (which will be in terms of random multiplicative functions, which we transfer to in the proof). Equipped with Propositions \ref{chaining}, \ref{bad_chaining_moment} and \ref{good_moment_bound} and , we can prove our theorems on the moment of moments. We finally note that it is sufficient to establish our bounds on the range $r \in[2/3,1]$ as then it suffices to apply Holder's inequality to obtain the result for smaller non negative $r$.
\begin{proof}[Proof of Theorems \ref{pf_moments_constant} and \ref{pf_moments_meso}, assuming Propositions \ref{bad_chaining_moment} and\ref{good_moment_bound}]
    This proof follows in three parts: approximate $|L(1/2+ih,\chi)|^2$ as a product of three Dirichlet series, extract the conditioned part and then handle the log-correlated part as in \cite{harper2019partition}. The first step is not too difficult. An application of Lemma \ref{Dirichlet_AFE} (and dropping the conditioning on $Y_{P_\theta}(0)$) gives
    \begin{align*}
        &\mathbb{E}_W^{\text{char}}\big(\int_{|h| \leq \log^{\theta}(q)/2}|L(1/2+ih,\chi)|^2dh\big)^r 
        \ll \mathbb{E}_W^{\text{char}}\big(\int_{|h| \leq \log^\theta(q)/2}|\sum_{n \leq \sqrt{q/2\pi}}\frac{\chi(n)}{n^{1/2+ih}}|^2dh\big)^r + q^{-r/2}.
    \end{align*}
    Applying the triangle inequality to the first term in the last line provides an upper bound of
    \begin{align*}
        &\mathbb{E}^{\text{char}}\big(\int_{|h| \leq \log^{\theta}(q)/2}|\sum_{n \leq \sqrt{q/2\pi}}\frac{\chi(n)}{n^{1/2+ih}}-\sum_{\substack{m \leq q^{\varepsilon_\theta} \\ P^+(m) \leq P_\theta}}\frac{\chi(m)}{m^{1/2+ih}}\sum_{\substack{n \leq q^{1/2-2\varepsilon_\theta} \\ P^-(n) > P_\theta}}\frac{\chi(n)}{n^{1/2+ih}}|^2dh\big)^r \\ 
        +& \mathbb{E}_W^{\text{char}}\big(\int_{|h| \leq \log^\theta(q)/2}|\sum_{\substack{m \leq q^{\varepsilon_\theta} \\ P^+(m) \leq P_\theta}}\frac{\chi(m)}{m^{1/2+ih}}\sum_{\substack{n \leq q^{1/2-2\varepsilon_\theta} \\ P^-(n) > P_\theta}}\frac{\chi(n)}{n^{1/2+ih}}|^2dh\big)^r
    \end{align*}
    On the first term, we can ignore the conditioning on the value of $Y_{P_\theta}(0)$, and apply Holder's inequality on the sum of Dirichlet characters to show that it is
    $\ll \big(\log^{1+2\theta}(q)/(\log\log(q))^2\big)^r$,
    which is more than acceptable for our theorems using Lemma \ref{factored_form}. We repeat this to access our product of three Dirichlet polynomials. We see 
    \begin{align*}
        &\mathbb{E}_W^{\text{char}}\big(\int_{|h| \leq \log^\theta(q)/2}|\sum_{\substack{m \leq q^{\varepsilon_\theta} \\ P^+(m) \leq P_\theta}}\frac{\chi(m)}{m^{1/2+ih}}\sum_{\substack{n \leq q^{1/2-2\varepsilon_\theta} \\ P^-(n) > P_\theta}}\frac{\chi(n)}{n^{1/2+ih}}|^2dh\big)^r \\
        \ll& \mathbb{E}_W^{\text{char}}\big(\int_{|h| \leq \frac{\log^\theta(q)}{2}}|\sum_{\substack{m \leq q^{\varepsilon_\theta} \\ P^+(m) \leq P_\theta}}\frac{\chi(m)}{m^{1/2+ih}}|^2|\sum_{\substack{n \leq q^{1/2-2\varepsilon_\theta} \\ P^{-}(n) >P_\theta}}\frac{\chi(n)}{n^{\sigma+ih}}-\sum_{\substack{k \leq q^{\varepsilon} \\ p|k \Rightarrow P_\theta < p \leq P}}\frac{\chi(k)}{k^{1/2+ih}}\sum_{\substack{n \leq q^{1/2-2\varepsilon} \\ P^{-}(n) > P}}\frac{\chi(n)}{n^{\sigma+ih}}|^2dh\big)^r \\
        +& \mathbb{E}_W^{\text{char}}\big(\int_{|h| \leq \log^\theta(q)/2}|\sum_{\substack{m \leq q^{\varepsilon_\theta} \\ P^+(m) \leq P_\theta}}\frac{\chi(m)}{m^{1/2+ih}}\sum_{\substack{k \leq q^{\varepsilon} \\ p|k \Rightarrow P_\theta < p \leq P}}\frac{\chi(k)}{k^{1/2+ih}}\sum_{\substack{n \leq q^{1/2-2\varepsilon} \\ P^{-}(n) > P}}\frac{\chi(n)}{n^{\sigma+ih}}|^2dh\big)^r.
    \end{align*}
    The first term requires us to apply Holder's inequality so we can take the expectation inside the integral, at which point we use Lemma \ref{partition_result} to approximate the conditioning on the small primes by a smooth partition of unity $g$, and then we transfer over to a statement about random multiplicative functions using Lemma \ref{randomisation}. Since we are now in the random multiplicative function setting, we can expand out the definition of $g \leq \mathbf{1}_{Y_{P_\theta}(0) \in [W,W+1]}+\delta_m$ for $\delta_m = \frac{1}{(\log\log(q))^2}$ say and notice that $\mathbf{1}_{Y_{P_\theta}(0) \in [W,W+1]}$ only affects the sum over the very small primes. This allows us to use the independence of the $f(p)$ to note that
    \begin{equation*}\begin{split}
        &\mathbb{E}_W\big(|\sum_{\substack{m \leq q^{\varepsilon_\theta} \\ P^+(m) \leq P_\theta}}\frac{f(m)}{m^{1/2+ih}}|^2|\sum_{\substack{n \leq q^{1/2-2\varepsilon_\theta} \\ P^{-}(n) >P_\theta}}\frac{f(n)}{n^{\sigma+ih}}-\sum_{\substack{k \leq q^{\varepsilon} \\ p|k \Rightarrow P_\theta < p \leq P}}\frac{f(k)}{k^{1/2+ih}}\sum_{\substack{n \leq q^{1/2-2\varepsilon} \\ P^{-}(n) > P}}\frac{f(n)}{n^{\sigma+ih}}|^2\big) \\
        =& \mathbb{E}_W\big(|\sum_{\substack{m \leq q^{\varepsilon_\theta} \\ P^+(m) \leq P_\theta}}\frac{f(m)}{m^{1/2+ih}}|^2\big)\mathbb{E}\big(|\sum_{\substack{n \leq q^{1/2-2\varepsilon_\theta} \\ P^{-}(n) >P_\theta}}\frac{f(n)}{n^{\sigma+ih}}-\sum_{\substack{k \leq q^{\varepsilon} \\ p|k \Rightarrow P_\theta < p \leq P}}\frac{f(k)}{k^{1/2+ih}}\sum_{\substack{n \leq q^{1/2-2\varepsilon} \\ P^{-}(n) > P}}\frac{f(n)}{n^{\sigma+ih}}|^2\big).
    \end{split}     
    \end{equation*}
    Then we can apply Proposition \ref{chaining} and Lemma \ref{logcor_factored_form} to show that
    $$\ll \big(e^{2W}\log^{1+2\theta}(q)/(\log\log(q))^2\big)^r.$$
    Now for $1-\frac{1}{\sqrt{\log\log(q)}} \leq r \leq 1$, we can apply Holder's inequality to the sum over Dirichlet characters, and we find this is
    \begin{align*}
        &\mathbb{E}_W^{\text{char}}\big(\int_{|h| \leq \log^\theta(q)/2}|\sum_{\substack{m \leq q^{\varepsilon_\theta} \\ P^+(m) \leq P_\theta}}\frac{\chi(m)}{m^{1/2+ih}}\sum_{\substack{k \leq q^{\varepsilon} \\ p|k \Rightarrow P_\theta < p \leq P}}\frac{\chi(k)}{k^{1/2+ih}}\sum_{\substack{n \leq q^{1/2-2\varepsilon} \\ P^{-}(n) > P}}\frac{\chi(n)}{n^{\sigma+ih}}|^2dh\big)^r \\
        \ll&\big(\int_{|h| \leq \log^\theta(q)/2}\mathbb{E}_W^{\text{char}}|\sum_{\substack{m \leq q^{\varepsilon_\theta} \\ P^+(m) \leq P_\theta}}\frac{\chi(m)}{m^{1/2+ih}}\sum_{\substack{k \leq q^{\varepsilon} \\ p|k \Rightarrow P_\theta < p \leq P}}\frac{\chi(k)}{k^{1/2+ih}}\sum_{\substack{n \leq q^{1/2-2\varepsilon} \\ P^{-}(n) > P}}\frac{\chi(n)}{n^{\sigma+ih}}|^2dh\big)^r \\
        \ll& \big(\int_{|h| \leq \log^\theta(q)/2}\mathbb{E}_W|\sum_{\substack{m \leq q^{\varepsilon_\theta} \\ P^+(m) \leq P_\theta}}\frac{f(m)}{m^{1/2+ih}}|^2\mathbb{E}|\sum_{\substack{k \leq q^{\varepsilon} \\ p|k \Rightarrow P_\theta < p \leq P}}\frac{f(k)}{k^{1/2+ih}}\sum_{\substack{n \leq q^{1/2-2\varepsilon} \\ P^{-}(n) > P}}\frac{f(n)}{n^{\sigma+ih}}|^2dh\big)^r \\
        \ll& \big(\log^{1+\theta}(q)\int_{|h| \leq \log^\theta(q)/2}\mathbb{E}_W|\sum_{\substack{m \leq q^{\varepsilon_\theta} \\ P^+(m) \leq P_\theta}}\frac{f(m)}{m^{1/2+ih}}|^2dh\big)^r. 
    \end{align*}
    This an appropriate form to apply Proposition \ref{chaining}, and the above then is
    $\ll\bigg(e^{2W}\log^{1+2\theta}(q)\bigg)^r$,
    as required.\par
    Now, for the remaining values of $r \in [2/3,1-\frac{1}{\sqrt{\log\log(q)}}]$, we use an argument first used in \cite{aharper201} which follows by multiple applications of Holder's inequality. For the remaining $r$, we set
    \begin{align*}
        R(\alpha) := \sup_{1-2\alpha \leq r \leq 1-\alpha}\bigg[&\mathbb{E}_W^{\text{char}}\big(\int_{|h|\leq \log^\theta(q)/2}|\sum_{\substack{m \leq q^{\varepsilon_\theta} \\ P^+(m) \leq P_\theta}}\frac{\chi(m)}{m^{1/2+ih}}\sum_{\substack{k \leq q^{\varepsilon} \\ p|k \Rightarrow P_\theta < p \leq P}}\frac{\chi(k)}{k^{1/2+ih}}\sum_{\substack{n \leq q^{1/2-2\varepsilon} \\ P^{-}(n) > P}}\frac{\chi(n)}{n^{\sigma+ih}}|^2dh\big)^r \\
        \cdot& \big(\frac{(1-r)\sqrt{\log\log(q)}}{\log^{1+\theta}(q)}\big)^r\bigg].
    \end{align*}
    We split this according to whether $\mathcal{G}_\chi$ holds or not. On this event holding, we apply Holder's inequality, approximate the indicator $\mathbf{1}_{Y_{P_\theta}(0) \cap \mathcal{G}_\chi}$ by a smooth function $g$ provided to us by Lemma \ref{partition_result} and then can pass to random multiplicative functions by Lemma \ref{randomisation}. Then we can upper bound $g$ from Lemma \ref{partition_result} (using $\delta=1/(\log\log(P))^2$ say), and apply independence of conditioning on the smooth and log-correlated part of the sums with orthogonality so we can finally apply Proposition \ref{good_moment_bound}. Putting all these steps together yields
    \begin{align*}
        &\mathbb{E}_W^{\text{char}}\big(\mathbf{1}_{\mathcal{G}_\chi}\int_{|h|\leq \log^\theta(q)/2}|\sum_{\substack{m \leq q^{\varepsilon_\theta} \\ P^+(m) \leq P_\theta}}\frac{\chi(m)}{m^{1/2+ih}}\sum_{\substack{k \leq q^{\varepsilon} \\ p|k \Rightarrow P_\theta < p \leq P}}\frac{\chi(k)}{k^{1/2+ih}}\sum_{\substack{n \leq q^{1/2-2\varepsilon} \\ P^{-}(n) > P}}\frac{\chi(n)}{n^{\sigma+ih}}|^2dh\big)^r \\
        \ll &\big(\int_{|h|\leq \log^\theta(q)/2}\mathbb{E}_W|\sum_{\substack{m \leq q^{\varepsilon_\theta} \\ P^+(m) \leq P_\theta}}\frac{f(m)}{m^{1/2+ih}}|^2\mathbb{E}\big(\mathbf{1}_{\mathcal{G}(h)}|\sum_{\substack{k \leq q^{\varepsilon} \\ p|k \Rightarrow P_\theta < p \leq P}}\frac{f(k)}{k^{1/2+ih}}\sum_{\substack{n \leq q^{1/2-2\varepsilon} \\ P^{-}(n) > P}}\frac{f(n)}{n^{\sigma+ih}}|^2\big)dh\big)^r \\
        \ll&\big(C\log^{1+\theta}(q)\min\{1,\frac{1}{(1-r)\sqrt{\log\log(P)}}\}\int_{|h| \leq \log^\theta(q)/2}\mathbb{E}_W|\sum_{\substack{m \leq q^{\varepsilon_\theta} \\ P^+(m) \leq P_\theta}}\frac{f(m)}{m^{1/2+ih}}|^2dh\big)^r
    \end{align*}
    where we achieved the last line by applying Proposition \ref{good_moment_bound}. Once again, we look to Proposition \ref{chaining}, so we obtain
    \begin{align*}
        &\mathbb{E}_W^{\text{char}}\big(\mathbf{1}_{\mathcal{G}_\chi}\int_{|h|\leq \log^\theta(q)/2}|\sum_{\substack{m \leq q^{\varepsilon_\theta} \\ P^+(m) \leq P_\theta}}\frac{\chi(m)}{m^{1/2+ih}}\sum_{\substack{k \leq q^{\varepsilon} \\ p|k \Rightarrow P_\theta < p \leq P}}\frac{\chi(k)}{k^{1/2+ih}}\sum_{\substack{n \leq q^{1/2-2\varepsilon} \\ P^{-}(n) > P}}\frac{\chi(n)}{n^{\sigma+ih}}|^2dh\big)^r \\
        \ll& \big(Ce^{2(A+W)}\log^{1+\theta}(q)\min\{1,\frac{1}{(1-r)\sqrt{\log\log(q)}}\}\big)^r \\ \ll& C\big(e^{2W}\log^{1+\theta}(q)\min\{1,\frac{1}{(1-r)\sqrt{\log\log(q)}}\}\big)^{1-\alpha}.
    \end{align*}
    Now, when we have $\mathcal{G}_\chi$ does not hold, then we set $r' = (1+r)/2$, so we have $1-\alpha \leq r' \leq 1-\alpha/2$. Then we apply Holder's inequality with exponents $r'/(r'-r)$ and $r'/r$ to get
    \begin{align*}
        &\mathbb{E}_W^{\text{char}}\big(\mathbf{1}_{\mathcal{G}_\chi \text{ fails}}\int_{|h|\leq \log^\theta(q)/2}|\sum_{\substack{m \leq q^{\varepsilon_\theta} \\ P^+(m) \leq P_\theta}}\frac{\chi(m)}{m^{1/2+ih}}\sum_{\substack{k \leq q^{\varepsilon} \\ p|k \Rightarrow P_\theta < p \leq P}}\frac{\chi(k)}{k^{1/2+ih}}\sum_{\substack{n \leq q^{1/2-2\varepsilon} \\ P^{-}(n) > P}}\frac{\chi(n)}{n^{\sigma+ih}}|^2dh\big)^r \\
        \ll& \big(\mathbb{E}^{\text{char}}\mathbf{1}_{\mathcal{G}_\chi \text{ fails}}\big)^{(r'-r)/r'} \\
        \cdot& \big(\mathbb{E}_W^{\text{char}}\big(\int_{|h|\leq \log^\theta(q)/2}|\sum_{\substack{m \leq q^{\varepsilon_\theta} \\ P^+(m) \leq P_\theta}}\frac{\chi(m)}{m^{1/2+ih}}\sum_{\substack{k \leq q^{\varepsilon} \\ p|k \Rightarrow P_\theta < p \leq P}}\frac{\chi(k)}{k^{1/2+ih}}\sum_{\substack{n \leq q^{1/2-2\varepsilon} \\ P^{-}(n) > P}}\frac{\chi(n)}{n^{\sigma+ih}}|^2dh\big)^{r'}\big)^{r/r'}.
    \end{align*}
    The contribution from the first term is $\ll e^{-C/2}$ using that $(r'-r)/r' \geq \alpha/2$ and from Proposition \ref{bad_chaining_moment}, we have $\mathbb{E}^{\text{char}}\big(\mathbf{1}_{\mathcal{G}_\chi \text{ fails}}\big) \ll e^{-C/\alpha}$. In particular, we have
    $$R(\alpha) \ll Ce^{2W} + e^{-C/2}(e^{2W}+R(\alpha/2)) \ll Ce^{2W} + e^{-C/2}R(\alpha/2).$$
    Successive applications of the above replacing $\alpha$ with $\alpha/2, \alpha/4,...$ gain the recursive bound
    $$R(\alpha) \ll e^{2W}+R(1/\sqrt{\log\log(q)}).$$
    From our previous discussion on the regime $1-\frac{1}{\sqrt{\log\log(q)}} \leq r \leq 1$, we know $R(1/\sqrt{\log\log(q)}) \ll e^{2W}$, which allows us to conclude, once we have made a suitable choice of $P$ and $\varepsilon$. We can choose $P=q^{1/(\log\log(q))^8}$ and $\varepsilon = \frac{1}{(\log\log(P))^2}+\varepsilon_\theta$, which are more than sufficient for our purposes.
\end{proof}
\subsection{Proof of the Proposition \ref{bad_chaining_moment}  and \ref{good_moment_bound}} As Proposition \ref{bad_chaining_moment} is slightly simpler, we begin there.
\begin{proof}[Proof of Proposition \ref{bad_chaining_moment}]
    For convenience, it is much easier to write $\mathcal{G}_\chi$ as a statement of a sum of logarithms of partial Euler products. From the union bound over all the $h(j)$, we can upper bound $\mathbb{P}_{q}(\mathcal{G}_\chi \text{ fails})$ by
    \begin{equation}
        \sum_{j=0}^{\lfloor\log(P_q(\theta))\rfloor -B-1}\sum_{h(j)}\mathbb{P}_q\big(\big|\sum_{\ell=j}^{\lfloor\log(P_q(\theta))\rfloor -B-1}\log(|I_{\ell,\chi}(h(\ell))|)\big|>\log(P_q(\theta)) -j + g(j)\big) \label{union_fail}
    \end{equation}
    for $P_q(\theta) = \log(P)\log^\theta(q)$ (defined for brevity). The second sum is over all the possible positions of the $h(j)$ over the interval $|h| \leq \log^\theta(q)$, where there are \\ $\ll (P_q(\theta)/e^j)(\log(\log(P)/e^j)$ such points, and once we know $h(j)$ then the points $h(\ell)$ are uniquely determined for $\ell \geq j$. Taylor expanding $\log(|I_\ell,\chi(h)|)$ reveals the sum over the $\ell$ is equal to
    \begin{align*}
        -\Re&\big(\sum_{\ell=j}^{\lfloor\log(P_q(\theta)\rfloor - B -1}\sum_{P^{e^{-(\ell+1)}} < p \leq P^{e^{-\ell}}}\log(1-\frac{\chi(p)}{p^{1/2+ih(\ell)}})\big) \\
        =\Re&\big(\sum_{\ell=j}^{\lfloor\log(P_q(\theta)\rfloor - B -1}\sum_{\substack{P^{e^{-(\ell+1)}} < p \leq P^{e^{-\ell}} \\ v=1,2}}\frac{\chi(p)}{p^{1/2+ih(\ell)}}+\frac{(\chi(p))^2}{p^{1+2ih(\ell)}}+O(p^{-3/2})\big).
    \end{align*}
    Consequently, if $\mathcal{G}_\chi$ fails, then we must have
    $$\big|\sum_{\ell=j}^{\lfloor\log(P_q(\theta)\rfloor - B -1}\sum_{\substack{P^{e^{-(\ell+1)}} < p \leq P^{e^{-\ell}} \\ v=1,2}}\big(\frac{\chi(p)}{p^{1/2+ih(\ell)}}+\frac{(\chi(p))^2}{p^{1+2ih(\ell)}}\big)\big| > \log(P_q(\theta)) + g(j)-j+O(1).$$
    An application of Markov's inequality to the power $$r=2\big\lfloor\frac{(\log(P_q(\theta))+g(j)-j+O(1))^2}{\log(P_q(\theta))-g(j)+O(1)}\big\rfloor$$
    combined with Stirling approximation yields
    \begin{align*}
        \sqrt{r}\exp\big(-\big\lfloor\frac{(\log(P_q(\theta))+g(j)-j+O(1))^2}{\log(P_q(\theta))-g(j)+O(1)}\big\rfloor\big) \ll q\sqrt{r}\frac{e^j}{P_q(\theta)}e^{-2g(j)-\frac{(g(j)+O(1))^2}{\log(P_q(\theta)-j+O(1)}} \numberthis \label{excep_set}
    \end{align*}
    where this has been achieved by an application of Lemma \ref{even_moment}. One deduces that $r$ is close to optimal by minimising taking the $2r$-th moment in Markov's inequality and balancing this with the Stirling approximation. Now, we split into the cases dependent on the size of $2C\min\{\sqrt{\log(P_q(\theta)},\frac{1}{1-q}\}$. Now if $C\min\{\sqrt{\log(P_q(\theta))},\frac{1}{1-q}\} < \log(P_q(\theta))-j$, then $r \asymp \log(P_q(\theta))-j$ clearly, and we obtain the bound $$\ll \frac{e^j}{P_q(\theta)}\frac{1}{\log^{7/2}(P_q(\theta)/e^j)}e^{-2C\min\{\sqrt{\log(P_q(\theta))},\frac{1}{1-q}\}}.$$ Otherwise, $C\min\{\sqrt{\log(P_q(\theta))},\frac{1}{1-q}\} \geq \log(P_q(\theta))-j$, then we have that one can bound equation (\ref{excep_set}) by $$\frac{e^j}{P_q(\theta)}\frac{1}{\log^{5}(P_q(\theta)/e^j)}e^{-3C\min\{\sqrt{\log(P_q(\theta)},\frac{1}{1-q}\}}.$$
    Regardless, we see that equation (\ref{excep_set}) is $$\ll q\frac{e^j}{P_q(\theta)}\frac{1}{\log^{7/2}(P_q(\theta)/e^j)}e^{-2C\min\{\sqrt{\log(P_q(\theta))},\frac{1}{1-q}\}}.$$. Putting this back into equation (\ref{union_fail}), we see
    \begin{align*}
        \mathbb{P}_q(\mathcal{G}_\chi \text{ fails}) \ll& \sum_{j=0}^{\lfloor\log(P_q(\theta))\rfloor -B-1}\sum_{h(j)}\frac{e^j}{P_q(\theta)}\frac{1}{\log^{7/2}(P_q(\theta)/e^j)}e^{-2C\min\{\sqrt{\log(P_q(\theta))},\frac{1}{1-q}\}} \\
        \ll& \sum_{j=0}^{\lfloor\log(P_q(\theta))\rfloor -B-1}\frac{1}{\log^{5/2}(P_q(\theta)/e^j)}e^{-2C\min\{\sqrt{\log(P_q(\theta))},\frac{1}{1-q}\}} \\
        \ll& e^{-2C\min\{\sqrt{\log(P_q(\theta))},\frac{1}{1-q}\}},
    \end{align*}
    using that $\log(P_q(\theta)/e^j) \asymp \log(\log(P)/e^j)$ as $\log^{1-\delta}(q) = o(\log(P))$ for any $\delta > 0$ with our choice of $P$.
\end{proof}
In order to prove Proposition \ref{good_moment_bound}, we will need another lemma, which can be imported from \cite{harper2019partition} with very minimal tweaks required (they are the same as the identifications used to prove Lemma \ref{logcor_factored_form} from Lemma \ref{factored_form}). We let $\mathcal{G}_{\text{rand}}(h)$ the event where $\mathcal{G}_{\text{rand}}$ occurs for a specific value $h$ (rather than occurring for each $h$ in the range).
\begin{lemma}
    Let $f$ be a Steinhaus multiplicative function and fix $\theta \in (-1/2,0]$. Then for $P_\theta < P \leq \sqrt{q}$ and $\varepsilon_\theta < \varepsilon \leq 1/10$, $2/3 \leq r \leq 1$ and $|h| \leq \log^\theta(q)/2$, we have
    \begin{align*}
        &\mathbb{E}\big(\mathbf{1}_{\mathcal{G}_{\text{rand}}(h)}|\sum_{\substack{m \leq q^\varepsilon \\ p|m \Rightarrow P_\theta < p \leq P}}\frac{f(m)}{m^{1/2+ih}}|^2|\sum_{\substack{n \leq q^{1/2-2\varepsilon} \\ P^-(n) > P}}\frac{f(n)}{n^{1/2+ih}}|^2\big) \\ &\ll \log^{1+\theta}(q)\big(C\min\big\{1,\frac{1}{(1-r)\sqrt{\log\log(P)}}\big\}+e^{-\frac{\varepsilon\log(q)}{\log(P)}}\big).
    \end{align*} \label{log_cor_barrier_moment}
\end{lemma}
While we will not prove the above, we will outline how the proof works. Since we are looking at a Dirichlet polynomial involving $P$ smooth numbers (with a barrier applied), one can relate these to a genuine truncated Euler product which also has the barrier present. Once we are working with random Euler products, then the barrier will control the growth of the random Euler product in the same way as shown in \cite{aharper201} in Section 5 (in the proof of Key Proposition 1). The only change required in our situation is that we are also restricting to the $P_\theta$ rough integers as well, which saves a factor of $\log(P_\theta)$ in both the standard expectation and the expectation with the barrier event applied. With this in mind, we are ready to prove Proposition \ref{good_moment_bound}.
\begin{proof}[Proof of Proposition \ref{good_moment_bound}]
    This statement becomes much simpler to prove given that we have Lemma \ref{randomisation} to handle the various error terms associated with introducing the partition of unity. We begin with an application of Holder's inequality over the sum of Dirichlet characters, which shows
    \begin{align*}
        &\mathbb{E}^{\text{char}}\big(\int_{|h| \leq \log^\theta(q)/2} \mathbf{1}_{\mathcal{G}_\chi}|\sum_{\substack{m \leq q^\varepsilon \\ p|m \Rightarrow P_\theta < p \leq P}}\frac{\chi(m)}{m^{1/2+ih}}|^2|\sum_{\substack{n \leq q^{1/2-2\varepsilon} \\ P^{-}(n) > P}}\frac{\chi(n)}{n^{1/2+ih}}|^2dh\big)^r \\
        \ll&\big(\int_{|h| \leq \log^\theta(q)/2} \mathbb{E}^{\text{char}}\big(\mathbf{1}_{\mathcal{G}_\chi}|\sum_{\substack{m \leq q^\varepsilon \\ p|m \Rightarrow P_\theta < p \leq P}}\frac{\chi(m)}{m^{1/2+ih}}|^2|\sum_{\substack{n \leq q^{1/2-2\varepsilon} \\ P^{-}(n) > P}}\frac{\chi(n)}{n^{1/2+ih}}|^2\big)dh\big)^r,
    \end{align*}
    so our goal will be to show
    \begin{align*}
        \mathbb{E}^{\text{char}}&\big(\mathbf{1}_{\mathcal{G}_\chi}|\sum_{\substack{m \leq q^\varepsilon \\ p|m \Rightarrow P_\theta < p \leq P}}\frac{\chi(m)}{m^{1/2+ih}}|^2|\sum_{\substack{n \leq q^{1/2-2\varepsilon} \\ P^{-}(n) > P}}\frac{\chi(n)}{n^{1/2+ih}}|^2dh\big) \\ \ll& \log^{1+\theta}(q)\min\big\{1,\frac{1}{(1-r)\sqrt{\log\log(P)}}\big\}.
    \end{align*}
    Like in the proof of Proposition \ref{chaining}, it is very convenient to bound $\mathbf{1}_{\mathcal{G}_\chi}(h)$ by $\mathbf{1}_{\mathcal{G}_\chi(h)}(h)$, where $\mathcal{G}_\chi(h)$ is the event such that for all $0 \leq j \leq \log(P_q(\theta)) - B -1$, we have
$$\big(\frac{P_q(\theta)}{e^j}e^{g(j)}\big)^{-1} \leq \prod_{\ell=j}^{\lfloor\log(P_q(\theta))\rfloor-B-1}I_{\ell,\chi}(h(\ell))\leq \frac{P_q(\theta)}{e^j}e^{g(j)}.$$
To deduce the parameters we will need to apply Lemma \ref{partition_result}, we recall that
$$\log\big(\prod_{\ell=j}^{\lfloor\log(P_q(\theta))\rfloor - B -1}|I_{\ell,\chi}(h)|\big) = \Re\big(\sum_{\ell=j}^{\lfloor\log(P_q(\theta)\rfloor - B -1}\sum_{\substack{P^{e^{-(\ell+1)}} < p \leq P^{e^{-\ell}}}}\sum_{v \geq 1}\frac{(\chi(p))^v}{vp^{v(1/2+ih(\ell))}}\big).$$
The terms corresponding to $v \geq 3$ are convergent and can be bounded by $8/3$. This then gives us our choice of $N_j$, which is
$$N_j = \log(P_\theta(q))-j+g(j)+8/3.$$
We will choose $\delta>0$ later. Then we can bound $\mathbf{1}_{\mathcal{G}_\chi(h)}(h)$ by
$$\prod_{j=0}^{\lfloor\log(P_q(\theta))\rfloor-B-1}\sum_{|i| \leq N_j}g_{j(i)}\big(\Re\big(\sum_{\ell=j}^{\lfloor\log(P_q(\theta)\rfloor - B -1}\sum_{\substack{P^{e^{-(\ell+1)}} < p \leq P^{e^{-\ell}} \\ v=1,2}}\frac{(\chi(p))^v}{vp^{v(1/2+ih(\ell))}}\big)\big).$$
Here we denote $g_{j(i)}$ as the partition of unity for the index $j$ on the interval $[i-N_j,i-N_j+1]$. Then assuming that we can apply Lemma \ref{randomisation}, we have
\begin{align*}
    &\sum_{|i| \leq N_j}\prod_{j=0}^{\lfloor\log(P_q(\theta))\rfloor-B-1}g_{j(i)}\big(\Re\big(\sum_{\ell=j}^{\lfloor\log(P_q(\theta)\rfloor - B -1}\sum_{\substack{P^{e^{-(\ell+1)}} < p \leq P^{e^{-\ell}} \\ v=1,2}}\frac{(\chi(p))^v}{vp^{v(1/2+ih(\ell))}}\big)\big) \\
    =& \sum_{|i| \leq N_j}\prod_{j=0}^{\lfloor\log(P_q(\theta))\rfloor-B-1}g_{j(i)}\big(\Re\big(\sum_{\ell=j}^{\lfloor\log(P_q(\theta)\rfloor - B -1}\sum_{\substack{P^{e^{-(\ell+1)}} < p \leq P^{e^{-\ell}} \\ v=1,2}}\frac{(f(p))^v}{vp^{v(1/2+ih(\ell))}}\big)\big)\\+&O(\frac{N_j\log(x)}{(\log(P_q(\theta))\log(P))^{\log(P_q(\theta))/\delta^2}}).
\end{align*}
This error term is more than acceptable for our theorem with our eventual choice of $\delta$ as $N_j$ is $\asymp \log\log(P)$, and we will be taking $\log(P)$ to the power of something larger than $\log\log(P)$, so this is sufficient. Then we can use Property $(b)$ of the $g$, so we can reverse the smoothing, and we find the above is
\begin{align*}
    &\mathbb{E}^{\text{char}}\big(\mathbf{1}_{\mathcal{G}_\chi}|\sum_{\substack{k \leq q^\varepsilon \\ p|k \Rightarrow P_\theta < p \leq P}}\frac{\chi(k))}{k^{1/2+ih}}|^2|\sum_{\substack{n \leq q^{1/2-2\varepsilon} \\ P^{-}(n) > P}}\frac{\chi(n)}{n^{1/2+ih}}|^2\big)dh \numberthis \label{randomed}\\ \ll& \mathbb{E}\big((1+\delta)^{\log(P_q(\theta))}(\mathbf{1}_{\mathcal{G}_{\text{rand}}(h)}(h)+\delta)|\sum_{\substack{k \leq q^\varepsilon \\ p|k \Rightarrow P_\theta < p \leq P}}\frac{f(k)}{k^{1/2+ih}}|^2|\sum_{\substack{n \leq q^{1/2-2\varepsilon} \\ P^{-}(n) > P}}\frac{f(n)}{n^{1/2+ih}}|^2\big).
\end{align*}
Choosing $\delta = \frac{1}{\log\log(P)}$ for example then shows that
$$(1+\delta)^{\log(P_q(\theta))}(\mathbf{1}_{\mathcal{G}_{\text{rand}}(h)}(h)+\delta) \ll \mathbf{1}_{\mathcal{G}_{\text{rand}}(h)}(h)+\delta.$$
An application of Lemma \ref{log_cor_barrier_moment} and standard random Euler product estimates then show that (\ref{randomed}) is
$$\ll \log^{1+\theta}(q)\big(\min\big\{1,\frac{1}{(1-r)\sqrt{\log\log(P)}}\big\}+\frac{1}{\log\log(P)}\big).$$
All that there is now is to check whether our choice of parameters were suitable for applying Lemma \ref{randomisation}. Indeed, we have that $$xP^{400(Y/\delta)^2\log(N\log(P))} \ll q^{1/2-\varepsilon}P^{1200(\log\log(P))^7} \ll q^{1/2-\varepsilon}q^{1/\log\log(q)} = o(q^{3/4})$$ for our choice of $P = q^{1/(\log\log(q))^8}$ which is again sufficient for application of Lemma \ref{randomisation}.
\end{proof}
\section{Proof of Theorem \ref{short_range_max}} \label{typical_max_section}
With the help of the two following propositions, we prove Theorem \ref{short_range_max}. To state these propositions, we need to define a suitable barrier event on the log correlated part of our Dirichlet polynomial. This will be very similar to the event $\mathcal{G}_\chi$ defined in Section \ref{pfm_proof}. To do this, we introduce the following set of discrete points which we use to approximate $h$ at various levels of precision. Let $\tilde{h}(-1) =h$ and for $0 \leq j \leq \log\log(P)-1$, define
$$\tilde{h}(j) := \max\big\{u \leq \tilde{h}(j-1):u=\frac{n}{(\log(P)/e^j)(\log\log(P))^3} \text{ for some } n \in \mathbb{Z}\big\}.$$
We also define $\tilde{h}(*)$ as the point $\frac{n}{\log(q)}$ which is closest to $h$. We set $I_{\ell,\chi}(h)$ and $I_{\ell,\text{rand}}(h)$ as they are defined in equations (\ref{char_increment}) and (\ref{random_increment}) respectively. Then, we let $\widetilde{\mathcal{G}}_\chi$ denote the event that for all $0 \leq j \leq \log\log(P)-B-1$
$$\big|\sum_{\ell =j}^{\lfloor\log\log(P)\rfloor-B-1}\log(I_{\ell,\chi}(\tilde{h}(\ell)))\big| \leq \log\log(P)-j+3\log\log\log(P) +U$$
for all $|h| \leq 1/2$ with $B$ some large constant and $U\geq 0$ a parameter. We let $\widetilde{\mathcal{G}}_\chi(h)$ denote the event that this occurs for a particular $h$ (rather than all of them). The minor change to look $I_{\ell,\chi}(h)$ rather than $|I_{\ell,\chi}(h)|$ will simplify our proof for proving a suitable bound on the number of $\chi$ such that $\widetilde{\mathcal{G}}_\chi$ fails. Then our proof of Theorem \ref{short_range_max} can be obtained with the two propositions below.
\begin{proposition}
    Uniformly for all $P \leq q^{1/(\log\log(q))^2}$ and $0 \leq U \leq 2\log\log(P)$,
    $$\mathbb{P}_q\big(\widetilde{\mathcal{G}}_\chi \text{ fails}\big) \ll e^{-2U}.$$ \label{max_event_fail}
\end{proposition}
\begin{proposition}
    Uniformly for all $P \leq q^{1/(\log\log(q))^8}, \ \frac{20\log(P)\log\log(P)}{\log(q)}< \varepsilon < 1/10, \ 0 \leq U \leq 2\log\log(P)$ and $V \geq e^{-U}$; and for all $|h| \leq 1/2$ and $|\sigma-1/2| \leq 1/\log(q)$, we have
    \begin{align*}
        &\mathbb{E}^{\text{char}}\big(\mathbf{1}_{\widetilde{\mathcal{G}}_\chi(h)}(h)\mathbf{1}_{\mathcal{H}_\chi(h;V) \text{ fails}}(h)\big|\sum_{\substack{m \leq q^\varepsilon \\ P^+(m) \leq P}}\frac{\chi(m)}{m^{1/2+ih}}\sum_{\substack{n \leq q^{1/2-2\varepsilon} \\ P^-(n) > P}}\frac{\chi(n)}{n^{\sigma+ih}}\big|^2\big) \\
        \ll& \log(q)\min\big\{1,\frac{U+\log\log\log(P)}{\sqrt{\log\log(P)}}\big\}\big(\min\big\{1,\frac{U+\log\log\log(P)+\log(V)}{\sqrt{\log\log(P)}}\big\}\big)^2. 
    \end{align*} \label{conditional_bound_max}
\end{proposition}
Note this proof is exactly identical to the work of Harper in \cite{harper2019partition}, only at points we choose to condense heavy technical theorems to outlines.
\begin{proof}[Proof of Theorem \ref{short_range_max} assuming Propositions \ref{max_event_fail} and \ref{conditional_bound_max}.]
    We set $P= q^{1/(\log\log(q))^8}$ and $\varepsilon = \frac{1}{(\log\log(q))^2}$. We also set the parameter $V=e^{-U}(\log\log(q))^6$. All of these are suitable for the application of our various propositions. \par
    For each $\chi$ modulo $q$, we define $h_\chi$ as the value of $h$ such that $\max_{|h| \leq1/2}|L(1/2+ih,\chi)|$ is obtained. From the statement of Theorem \ref{short_range_max}, we will need to prove
    \begin{equation}
        \mathbb{P}_q\big(|L(1/2+ih_\chi,\chi)|>\frac{e^{U}\log(q)}{(\log\log(q))^{3/4}}\big) \ll e^{-2U}(U+\log\log\log(q))(\log\log\log(q))^2. \label{target_max}
    \end{equation}
    The size of this set is larger than the set of $\chi$ where $\widetilde{\mathcal{G}}_\chi$ fails, so we may assume that the event $\widetilde{\mathcal{G}}_\chi$ holds in our calculations. This will be particularly helpful when the log correlated contribution to $L(1/2+ih,\chi)$ is large. \par
    Now, if $|L(1/2+ih_\chi),\chi)| \geq \frac{e^{U}\log(q)}{(\log\log(q))^{3/4}}$, then by the approximate functional equation of Dirichlet $L$-functions, we know
    $$\sum_{n \leq \sqrt{q/2\pi}}\frac{\chi(n)}{n^{1/2+ih_\chi}} \gg \frac{e^{U}\log(q)}{(\log\log(q))^{3/4}},$$
    given the error in this approximation is $O(q^{-1/4})$ (which is acceptable given the bound in equation (\ref{target_max})). This Dirichlet polynomial is a holomorphic function, so we can apply the Cauchy integral formula to rewrite this. We then see
    $$\sum_{n \leq \sqrt{q/2\pi}}\frac{\chi(n)}{n^{1/2+ih_\chi}} = \frac{1}{2\pi i}\big(\int_{A_\chi}+\int_{B_\chi}+\int_{C_\chi}+\int_{D_\chi}\big)\sum_{n \leq \sqrt{q/2\pi}}\frac{\chi(n)}{n^{s}}\frac{ds}{(s-1/2-ih_\chi)},$$
    where the union of the lines $A_\chi, B_\chi, C_\chi$ and $D_\chi$ form a square with corners $1/2+ \tilde{h_\chi}(*) \pm\frac{1 \pm i}{\log(q)}$. From the way $\tilde{h}(*)$ is defined, we have $|h_\chi-\tilde{h_\chi}(*)| \leq 1/\log(q)$, and further we can see that within the square defined by $A_\chi, B_\chi, C_\chi, D_\chi$, we have $|s-1/2-ih_\chi| \ll 1/\log(q)$. In light of this, we have
    $$\max_{\mathcal{C} \in \{A_\chi,B_\chi,C_\chi,D_\chi\}}\int_{\mathcal{C}}\sum_{n \leq \sqrt{q/2\pi}}\big|\frac{\chi(n)}{n^{s}}\big|d|s|  \gg \frac{e^{U}}{(\log\log(q))^{3/4}}.$$
    We investigate the vertical lines first. These can also be handled in an identical fashion, so here $\sigma$ can denote $1/2+1/\log(q)$ or $1/2-\log(q)$. After several applications of the triangle inequality, we see after a change of variables $s \mapsto s+1/2$
    \begin{align*}
        &\int_{|\tilde{h_\chi}(*)-h| \leq 1/\log(q)}\sum_{n \leq \sqrt{q/2\pi}}\big|\frac{\chi(n)}{n^{\sigma+ih}}\big|dh  \numberthis \label{core_triangle}\\ 
        \ll& \int_{|\tilde{h_\chi}(*)-h| \leq 1/\log(q)}\big|\sum_{n \leq \sqrt{q/2\pi}}\frac{\chi(n)}{n^{\sigma+ih}}-\sum_{\substack{m \leq q^{\varepsilon} \\ P^+(m) \leq P}}\frac{\chi(m)}{m^{1/2+ih}}\sum_{\substack{n \leq q^{1/2-2\varepsilon} \\ P^-(n) > P}}\frac{\chi(n)}{n^{\sigma+ih}}\big|dh \\
        +& \int_{|\tilde{h_\chi}(*)-h| \leq 1/\log(q)}|\sum_{\substack{m \leq q^{\varepsilon} \\  P^+(m) \leq P}}\frac{\chi(k)}{k^{1/2+ih}}\sum_{\substack{n \leq q^{1/2-2\varepsilon} \\ P^{-}(n) > P}}\frac{\chi(n)}{n^{\sigma+ih}}|dh.
    \end{align*}
    For the first expression, if it is $\gg \frac{e^{U}}{(\log\log(q))^{3/4}}$, then by applying Cauchy-Schwarz, we must have
    $$\frac{e^{2U}}{(\log\log(q))^{3/2}} \ll \frac{1}{\log(q)}\int_{|\tilde{h_\chi}(*)-h| \leq 1/\log(q)}\big|\sum_{n \leq \sqrt{q/2\pi}}\frac{\chi(n)}{n^{\sigma+ih}}-\sum_{\substack{m \leq q^{\varepsilon} \\ P^+(m) \leq P}}\frac{\chi(m)}{m^{1/2+ih}}\sum_{\substack{n \leq q^{1/2-2\varepsilon} \\ P^-(n) > P}}\frac{\chi(n)}{n^{\sigma+ih}}\big|^2dh.$$
    We then can extend the range of integration to $|h| \leq 1/2$ (so we now have no character dependency), and sum over all Dirichlet characters modulo $q$. The number of $\chi$ where this occurs is then
    $$\ll \frac{qe^{-2U}}{\sqrt{\log\log(q)}},$$
    which is certainly good enough. \par
    The final term in equation (\ref{core_triangle}) requires a more delicate procedure. By Holder's inequality, if the final integral is $\gg \frac{e^{U}}{(\log\log(q))^{3/4}}$, then using the exponents $r=4/3$ and $r'=4$, then we have either
    \begin{align*}
        \frac{e^{4U}}{(\log\log(q))^3} \ll  \frac{1}{\log^3(q)}\int_{|\tilde{h_\chi}(*)-h| \leq 1/\log(q)}\mathbf{1}_{\mathcal{H}_\chi(h;V)}(h)\big|\sum_{\substack{m \leq q^\varepsilon \\ P^+(m) \leq P}}\frac{\chi(k)}{k^{1/2+ih}}\sum_{\substack{n \leq q^{1/2-2\varepsilon} \\ P^{-}(n) > P}}\frac{\chi(n)}{n^{\sigma+ih}}\big|^4dh \numberthis \label{small_dirichlet}
    \end{align*}
    or
    \begin{align*}
        &\frac{e^{2U}}{(\log\log(q))^{3/2}} \ll \int_{|\tilde{h_\chi}(*)-h| \leq 1/\log(q)}\mathbf{1}_{\mathcal{H}_\chi(h;V) \text{ fails}}(h)\big|\sum_{\substack{m \leq q^\varepsilon \\ P^+(m) \leq P}}\frac{\chi(m)}{m^{1/2+ih}}\sum_{\substack{n \leq q^{1/2-2\varepsilon} \\ P^{-}(n) > P}}\frac{\chi(n)}{n^{\sigma+ih}}\big|^2dh. \numberthis \label{big_dirichlet}
    \end{align*}
    We begin with equation (\ref{small_dirichlet}). Like before, in the second equation we extend the range of integration in $h$ to make the quantity independent of $\chi$, and sum over the Dirichlet characters modulo $q$. The integrand in equation (\ref{small_dirichlet}) is then handled by Lemma \ref{conditional_fourth}. From this, we deduce that the number of characters modulo $q$ that satisfy equation (\ref{small_dirichlet}) are then
    $$\ll \frac{qe^{-4U}(\log\log(q))^3}{V^2} \ll \frac{qe^{-2U}}{\log\log(q)}.$$
    Finally, we handle equation (\ref{big_dirichlet}). This is where we can use the event $\widetilde{\mathcal{G}}_\chi(h)$. We again extend the range of integration in $h$, we have
    $$\frac{e^{2U}}{(\log\log(q))^{3/2}} \ll \frac{1}{\log(q)}\int_{|h| \leq 1/2}\mathbf{1}_{\mathcal{H}_\chi(h;V) \text{ fails}}(h)\big|\sum_{\substack{m \leq q^\varepsilon \\ P^+(m) \leq P}}\frac{\chi(m)}{m^{1/2+ih}}\sum_{\substack{n \leq q^{1/2-2\varepsilon} \\ P^{-}(n) > P}}\frac{\chi(n)}{n^{\sigma+ih}}\big|^2dh.$$
    Then the number of $\chi$ such that the above display holds, and that $\widetilde{\mathcal{G}}_\chi$ satisfies
    $$\frac{(\log\log(q))^{3/2}e^{-2U}}{\log(q)}\int_{|h| \leq 1/2}\mathbb{E}^{\text{char}}\big(\mathbf{1}_{\widetilde{\mathcal{G}}_\chi}(h)\mathbf{1}_{\mathcal{H}_\chi(h;V) \text{ fails}}(h)\big|\sum_{\substack{m \leq q^\varepsilon \\ P^+(m) \leq P}}\frac{\chi(m)}{m^{1/2+ih}}\sum_{\substack{n \leq q^{1/2-2\varepsilon} \\ P^{-}(n) > P}}\frac{\chi(n)}{n^{\sigma+ih}}\big|^2\big)dh.$$
    From this, we use $\mathbf{1}_{\widetilde{\mathcal{G}}_\chi}(h) \leq \mathbf{1}_{\widetilde{\mathcal{G}}_\chi(h)}(h)$, so the above is
    \begin{align*}
        \ll& q(\log\log(q))^{3/2}e^{-2U}\min\big\{1,\frac{U+\log\log\log(P)}{\sqrt{\log\log(P)}}\big\}\big(\min\big\{1,\frac{U+\log\log\log(P)+\log(V)}{\sqrt{\log\log(P)}}\big\}\big)^2 \\
        \ll& qe^{-2U}(U+\log\log\log(q))(\log\log\log(q))^2.
    \end{align*}
    Now, it only remains to handle the contribution of the horizontal sides. These will be of the form
    $$\int_{|\sigma-1/2| \leq1/\log(q)}\big|\sum_{n \leq \sqrt{\frac{q-|a|/\log(q)}{2\pi}}}\frac{\chi(n)}{n^{\sigma+i(\tilde{h_\chi}(*)+a/\log(q))}}\big|d\sigma$$
    where $a\in\{-1,+1\}$. Now, we can use the triangle inequality to obtain an upper bound for the above in the same style as equation \ref{core_triangle}. What is quite helpful is that the shift in $\sigma$ only affects the Dirichlet polynomial supported on rough numbers (and the shift in $\sigma$ is so small it makes very little difference anyway). What does need more work is how we make each expression independent of $\chi$. This is specifically why we introduced the new approximating points $\tilde{h_\chi}(*)$ as these are all of the form $n/\log(q)$. We can then sum over all of the possible positions of $\tilde{h_\chi}(*)$ over the interval $[-1/2,1/2]$, of which there are $\asymp \log(q)$ of these points. One then can proceed in an identical manner to the way we do for the vertical integrals, and obtain the same bounds (in fact it is slightly easier because we are only handling a shift in the real part and not in the imaginary part of the argument). This is why all of our preliminary lemmas and propositions only involved averaging over the characters and uniformly in $h$ so they can be applied in this situation.
\end{proof}
\subsection{Proof of Proposition \ref{max_event_fail} and \ref{conditional_bound_max}}
Like in the case of our previous theorems, it is much easier to prove Proposition \ref{max_event_fail} than Proposition \ref{conditional_bound_max}, so we start there.
\begin{proof}[Proof of Proposition \ref{max_event_fail}]
    Most of the proof has already been completed in Proposition \ref{bad_chaining_moment}. Notice that this time, we are interested in bounding absolute values of the sum of $\log(I_{\ell,\chi}(\tilde{h}(\ell)))$, and not the sums of the real part of $\log(I_{\ell,\chi}(\tilde{h}(\ell)))$. This is not a problem as we upper bounded the real part by their absolute values in the proof of Proposition \ref{bad_chaining_moment}.
    We want to find the probability such that
    \begin{equation}
        \big|\sum_{\ell=j}^{\lfloor\log\log(P)\rfloor - B -1}\sum_{\substack{P^{e^{-(\ell+1)}} < p \leq P^{e^{-\ell}} \\ v=1,2}}\frac{\chi(p)}{p^{1/2+ih(\ell)}}+\frac{(\chi(p))^2}{p^{1+2ih(\ell)}}\big| > \log\log(P) -j + 3\log\log\log(P)+U. \label{max_excep_set}
    \end{equation}
    An application of Markov's inequality to the power
    $$r=\big\lfloor\frac{(\log\log(P)-j+3\log\log\log(P)+U+O(1))^2}{\log\log(P)-j+O(1)}\big\rfloor$$
    shows that the probability that \ref{max_excep_set} occurs for is
    $$\ll \sqrt{r}\big(\frac{r(\log\log(P)-j+O(1))}{e(\log\log(P)-j+3\log\log\log(P)+U+O(1))^2}\big)^r.$$
    As $U \leq 2\log\log(P)$, then we certainly know that $r \leq (3\log\log(P))^2$, so this is bounded by
    $$\ll q\frac{\log\log(P)e^j}{\log(P)(\log\log(P))^6}e^{-2U} \ll q\frac{e^j}{\log(P)(\log\log(P))^{5}}e^{-2U}.$$
    Substituting back into the initial union bound we take over the $\tilde{h}(j)$ (recall there are $\log(P)e^{-j}(\log\log(P))^3$ such points) and over $j$, we then see the final bound will be $\ll e^{-2U}$ as required.
\end{proof}
We need some preparation in order to prove Proposition \ref{conditional_bound_max}. This is considerably more technically demanding than Proposition \ref{good_moment_bound} since we are looking for such large values which are rare. Due to the more stringent barrier, we need something that resembles an Euler product more closely to work with to prove our result.
\begin{lemma}
    Suppose $P \leq q^{\frac{1}{100\log\log(q)}}$, with $\frac{e^2\log(P)\log\log(P)}{\log(q)}<\varepsilon \leq \frac{1}{(\log\log(P))^2}$. For each prime $p$, define $\ell(p) \in \mathbb{N} \cup \{0\}$ for which $P^{e^{-(\ell+1)}} < p \leq P^{e^{-\ell}}$. Then uniformly for all $|h| \leq 1/2$ and $|\sigma-1/2| \leq 1/\log(q)$
    \begin{align*}
        \mathbb{E}^{\text{char}}&\bigg(\big|\sum_{\substack{m \leq q^{\varepsilon} \\ P^+(m) \leq P}}\frac{\chi(m)}{m^{1/2+ih}}-\sum_{v=0}^{\lfloor\frac{\varepsilon\log(q)}{\log(P)}\rfloor}\big(\sum_{p^j:p \leq P}\frac{\chi(p)^j}{jp^{j(1/2+i\tilde{h}(\ell(p)))}}\big)^v\big|^2\bigg) \ll \frac{\log(P)}{(\log\log(P))^2}.
    \end{align*} \label{partial_exp}
\end{lemma}
\begin{proof}
    Again, this goes through without change to the analogous result Lemma 6 in \cite{harper2019partition} (only that we are using a character version for our even moment estimates).
\end{proof}
The barrier event we define on the random side will account for suitable random versions of the events $\widetilde{\mathcal{G}}_{\chi}(h)$ and $\mathcal{H}_\chi(h;V)$. The event $\widetilde{\mathcal{G}}_{\text{rand}}(h)$ will denote when for all $0 \leq j \leq \log\log(P)-B-1$ that we have
$$ \big|\sum_{\ell = j}^{\lfloor\log\log(P)\rfloor-B-1}\log(|I_{\ell,\text{rand}}(\tilde{h}(\ell))|)\big| \leq \log\log(P) -j + 3\log\log\log(P)+U+19/3$$
as well as the stronger lower bound
$$\prod_{j=0}^{\lfloor\log\log(P)\rfloor-B-1}|I_{\ell,\text{rand}}(\tilde{h}(\ell))| \geq \frac{\log(P)}{VB_0e^{19/3}}$$
for some absolute large constant $B_0$ (coming from the probability results in Section \ref{prob_heur}). One way that we might hope to obtain a stronger result here is to allow for the barriers to be curved with $j$ (which is done in \cite{abr20}). The barriers used in their work follow are far tighter than the ones used here as well (relying similar results to Lemma \ref{ballot} which allow for the slowly growing $h(j)$ function). We then need a suitable estimate for how our log-correlated Dirichlet polynomial behaves under this event. Again, this will follow in an identical manner to that of Lemma 7 in \cite{harper2019partition}. However, this is rather tricky, so we will expand upon why this is the case.
\begin{lemma}
    Let $f(n)$ be a Steinhaus multiplicative function. Uniformly for all large $P \leq \sqrt{q}$ and $0 < \varepsilon \leq 1/10$ with $0 \leq U \leq 2\log\log(P)$ and $V \geq e^{-U}$, then for all $|h| \leq 1/2$ and $|\sigma-1/2| \leq 1/\log(q)$,
    \begin{align*}
        &\mathbb{E}\big(\mathbf{1}_{\widetilde{\mathcal{G}}_{\text{rand}}(h)}(h)|\sum_{\substack{m \leq q^\varepsilon \\ P^+(m) \leq P}}\frac{f(m)}{m^{1/2+ih}}|^2|\sum_{\substack{n \leq q^{1/2-2\varepsilon} \\ P^-(n) > P}}\frac{f(n)}{n^{1/2+ih}}|^2\big)\\ \ll& \log(q)\big(\min\big\{1,\frac{U+\log\log\log(P)}{\sqrt{\log\log(P)}}\big\}\big(\min\big\{1,\frac{U+\log\log\log(P)+\log(V)}{\sqrt{\log\log(P)}}\big\}\big)^2+e^{-\frac{\varepsilon\log(q)}{\log(P)}}\big) 
    \end{align*} \label{log_cor_max_moment}
\end{lemma}
Unsurprisingly, this reduces to a bounding $\prod_{p \leq P}|1-\frac{f(p)}{p^{1/2+ih}}|^{-2}$ like in Lemma \ref{log_cor_barrier_moment}. Due to the small primes creating issues, a solution is then to condition on the values of Euler products involving the small primes, and then use the fact that these will be independent of the shorter Euler products not involving these small primes. When not handling the small primes, then the Gaussian comparison steps work in the same way as one sees in Lemma 4 in \cite{harper2019partition}, and one proceeds in the same way. If one is willing to permit the gaining of a factor $\ll \log_4(P)$ on the upper bound in Lemma \ref{log_cor_max_moment}, then these steps are unnecessary since we would not bother with conditioning on these small primes. The work in \cite{abr20} takes this idea further, since they completely do away with any kind of barrier condition on the Euler product increments over the small primes as their application of the twisted fourth moment alleviates issues concerning the smooth numbers.  \par
With this in mind, we can prove Proposition \ref{conditional_bound_max}.
\begin{proof}[Proof of Proposition \ref{conditional_bound_max}]
    To start with, we assume that $e^{-U} \leq V \leq e^{\sqrt{\log\log(P)}}$, otherwise we can ignore the indicator of the log-correlated part of the Dirichlet polynomial being large and proceed in the same way we did in Proposition \ref{good_moment_bound} (or we can proceed as we do here just ignoring all the manipulations to handle the size condition on the large Dirichlet polynomial). Expanding $\log(I_{\ell,\chi}(\tilde{h}(\ell))$, we see that $\widetilde{\mathcal{G}}_\chi(h)$ is operating on the quantity
    $$\big|\sum_{p^j:p \leq P}\frac{\chi(p)}{jp^{j(1/2+i\tilde{h}(\ell(p)))}}\big|=\big|\sum_{\ell =j}^{\lfloor\log\log(P)\rfloor-B-1}\log(I_{\ell,\chi}(\tilde{h}(\ell))) +O(1)\big| \leq 4\log\log(P).$$
    The $O(1)$ term is dependent on the large absolute constant $B$ from the probability results. Using another Taylor expansion on the exponential and noting $\frac{\varepsilon\log(q)}{\log(P)} > 20\log\log(P)$ (it is $\gg (\log\log(q))^6$ when we apply this to Theorem \ref{short_range_max}), we have
\begin{align*}
    \sum_{v=0}^{\lfloor\frac{\varepsilon\log(q)}{\log(P)}\rfloor}\frac{1}{(v!)}\big(\sum_{\substack{p^j:p \leq P}}\frac{\chi(p)^j}{jp^{j(1/2+i\tilde{h}(\ell(p))}}\big)^v =& \exp\big(\sum_{\substack{p^j:p \leq P}}\frac{(\chi(p))^j}{jp^{j(1/2+i\tilde{h}(\ell(p)))}}\big) + O\big(\frac{(4\log\log(P))^{\lceil\frac{\varepsilon\log(q)}{\log(P)}\rceil}}{(\left\lceil\frac{\varepsilon\log(q)}{\log(P)}\right\rceil)!}\big) \\
    =&e^{O(1)}\prod_{\ell=0}^{\lfloor\log\log(P)\rfloor-B-1}I_{\ell,\chi}(\tilde{h}(\ell)) + O(1).
\end{align*}
Again, the $e^{O(1)}$ term is dependent on fixed constant $B$ and the remainder term is bounded using the Stirling approximation. Why this matters is because we then can firmly connect $\widetilde{\mathcal{G}}_\chi(h)$ to our factored Dirichlet polynomial. More concretely, we then have there exists a constant $B_0$ dependent on $B$ such that
\begin{align*}
    &\mathbf{1}_{\widetilde{\mathcal{G}}_\chi(h)}(h)\mathbf{1}_{\mathcal{H}_\chi(h;V) \text{ fails}}(h)\big|\sum_{\substack{m \leq q^\varepsilon \\ P^+(m) \leq P}}\frac{\chi(m)}{m^{1/2+ih}}\sum_{\substack{n \leq q^{1/2-2\varepsilon} \\ P^-(n) > P}}\frac{\chi(n)}{n^{\sigma+ih}}\big|^2 \numberthis \label{event_translate}\\
    \leq& \mathbf{1}_{\widetilde{\mathcal{G}}_\chi}(h)\mathbf{1}_{\widetilde{\mathcal{H}}_\chi(h;\frac{V}{B_0}) \text{ fails}}\big|\sum_{\substack{m \leq q^\varepsilon \\ P^+(m)\leq P}}\frac{\chi(m)}{m^{1/2+ih}}\sum_{\substack{n \leq q^{1/2-2\varepsilon} \\ P^-(n) > P}}\frac{\chi(n)}{n^{\sigma+ih}}\big|^2 \\
    +& \big|\sum_{\substack{n \leq q^{1/2-2\varepsilon} \\ P^{-}(n) > P}}\frac{\chi(n)}{n^{\sigma+ih}}\big|^2\big|\sum_{\substack{k \leq q^{\varepsilon} \\ P^+(m) \leq P}}\frac{\chi(m)}{m^{1/2+ih}}-\sum_{v=0}^{\lfloor\frac{\varepsilon\log(q)}{\log(P)}\rfloor}\big(\sum_{p^j:p \leq P}\frac{\chi(p)^j}{jp^{j(1/2+i\tilde{h}(\ell(p))}}\big)^v\big|^2.
\end{align*}
This follows since if $|\sum_{\substack{m \leq q^\varepsilon \\ P^+(m) \leq P}}\frac{\chi(m)}{m^{1/2+ih}}|>\frac{\log(P)}{V}$ but $\prod_{\ell=0}^{\lfloor\log\log(P)\rfloor-B-1}I_{\ell,\chi}(\tilde{h}(\ell)) \leq \frac{\log(P)}{B_0V}$, then from Lemma \ref{partial_exp}, $|\sum_{\substack{m \leq q^\varepsilon \\ P^+(m) \leq P}}\frac{\chi(m)}{m^{1/2+ih}}|$ be comparable to the difference
$$\big|\sum_{\substack{m \leq q^\varepsilon \\ P^+(m) \leq P}}\frac{\chi(m)}{m^{1/2+ih}}-\sum_{v=0}^{\lfloor\frac{\varepsilon\log(q)}{\log(P)}\rfloor}\big(\sum_{p^j:p \leq P}\frac{\chi(p)^j}{jp^{j(1/2+i\tilde{h}(\ell(p)))}}\big)^v\big| \ll \frac{\sqrt{\log(P)}}{\log\log(P)}$$
by Cauchy-Schwarz (and this can't happen with our restriction on $V$). The second term in the inequality of the display (\ref{event_translate}) is $\ll \frac{\log(P)}{(\log\log(P))^2}$ once we sum over the Dirichlet characters modulo $q$. All that is left is the Dirichlet polynomial with barriers $\widetilde{\mathcal{G}}_\chi(h)$ and $\mathcal{H}_\chi(h;V)$ applied. We approximate these by a product of smooth functions provided to us by Lemma \ref{partition_result} in a similar manner to that done in Proposition \ref{good_moment_bound}, where we upper bound $\mathbf{1}_{\widetilde{\mathcal{G}}_\chi(h)}$ by the indicator for the event such that for all $0 \leq j \leq \lfloor\log\log(P)\rfloor-B-1$
$$\big|\sum_{\ell=j}^{\lfloor\log\log(P)\rfloor-B-1}I_{\ell,\chi}(h(\ell))\big|\leq \log\log(P)-j+3\log\log\log(P)+U.$$
The only difference will be that we will need to slightly adapt the parameters required for the $j=0$ case since we must incorporate the stronger lower bound condition. Assuming we satisfy the conditions of Lemma \ref{randomisation}, up to suitable error terms we have that first term after the inequality of equation \ref{event_translate} can be bounded above by
$$\ll (1+\delta)^{\log\log(P)}\mathbb{E}\big((\mathbf{1}_{\widetilde{\mathcal{G}}_{\text{rand}}(h)}(h)+\delta)|\sum_{\substack{m \leq q^\varepsilon \\ P^+(m)\leq P}}\frac{f(m)}{m^{1/2+ih}}|^2|\sum_{\substack{n \leq q^{1/2-2\varepsilon} \\ P^-(n) > P}}\frac{f(n)}{n^{1/2+ih}}|^2\big).$$
In order to achieve Proposition \ref{conditional_bound_max}, we can choose $\delta = \frac{1}{(\log\log(P))^2}$ so the second term in the expectation is bounded by $\ll \frac{\log(q)}{(\log\log(q))^2}$ and $(1+\delta)^{\log\log(P)} \ll 1$. The first term in the expectation is handled by Lemma \ref{log_cor_max_moment}, and all we have to check is that we can apply Lemma \ref{randomisation} with our choice of parameters, which are suitable in line with similar reasoning seen in the proof of Proposition \ref{good_moment_bound} for example.
\end{proof}
\printbibliography[title= References]
\end{document}
